% AI-written, AI cross-reviewed; no human sign-off; not a claim that RH is proved
% RH Notes blueprint source excerpts for Z4.2.13
% Verbatim excerpts in the notation of the parent blueprint. The labelled editions and source ranges appear below.
% See the unit YAML for the current proof status. Earlier-round passages retain their original wording.
\section*{Z4.2.13. Period, lattice and theta-vector laws of the tied base}

% Parent source SHA-256: 99a8461100a8be8176578b515f15d452958d95ac91f86554c8727ce0a98ee36e
% blueprint/ch_P_r14.tex lines 34 through 54; r14:period
\paragraph{Why the period cannot vanish.}\label{r14:period}
$B.L$ is a period pair whose lattice is $\mathrm{lamOf}(D,\mathfrak c,\delta)$
(\nolinkurl{Base.lat}, \path{KernelLattices/LateralTheory.lean:128--132}).
If $\delta=0$, each generator $\delta\phi_0(c)$ is zero, so its additive
closure is zero. The first period belongs to the lattice, whereas the first
vector of an independent period pair is nonzero. This contradiction uses
only the base, so it supplies the same nonvanishing hypothesis whenever
the lattice proofs below require it
(\nolinkurl{lamOf_zero}, \nolinkurl{base_linkage_nonempty},
\path{GlobalTheory/GL3/PFamily.lean:74--86}).

For the two statements without $B$ (\nolinkurl{TrueTheta.exists_coef_field}
and \nolinkurl{TrueTheta.exists_family_of_curve}), start with the unscaled
ideal lattice. A nonzero ideal gives an independent period pair there.
Its Weierstrass discriminant is nonzero, and
$\delta^{12}=g_2^3-27g_3^2$ for this lattice, hence $\delta\ne0$.
See \nolinkurl{CurveLattice.discr_ne_zero},
\nolinkurl{CurveLattice.exists_periodPair_lamOf}
(\path{FoundationProofs/ExistsPeriodPairIdealCurveLattice.lean:64--67,112--119}),
\nolinkurl{InertSign.delta_ne_zero}
(\path{WeilRepresentationVacuumSignInertPlaces.lean:240--248}).

% blueprint/ch_P_r14.tex lines 75 through 125; r14:lat-integral
\paragraph{Integrality and parity.}\phantomsection\label{r14:lat-integral}
Write $h(S,T)=\operatorname{Tr}_{E/K}\operatorname{tr}(T^*S)$.
The fixed host supplies that $I$ is its own dual for
$(a,b)\mapsto\operatorname{Tr}_{E/K}(a\bar b)$: its different is
$\mathfrak b\bar{\mathfrak b}$, so the trace dual of $\bar I$ is
$(\mathfrak b\bar{\mathfrak b})^{-1}(\bar I)^{-1}=I$.
For lattice homomorphisms $P,Q$, $Q^*$ sends $L'$ into $L$ by
self-duality. Thus $Q^*P$ preserves $L$, and its trace belongs to
$\mathcal O_E$: after localization at each prime, the full torsion-free
lattice is free and an integral matrix represents the endomorphism there.
For $a,b\in I$, absorb this integral trace into $a$. Trace duality gives
$h(aP,bQ)\in\mathcal O_K$, hence $h(M,M)\subset\mathcal O_K$ by addition.
The $\tau$ are precisely the extensions of $\phi_0|_K$ (the conjugate half
acts differently on $\sqrt{-d}$), so their sum computes the relative trace.
See \nolinkurl{HostDer.unimodular}, the host different factorization
(\path{TotientGaloisClasses/ConjEMemHostFIff.lean:159--162},
\path{LocalRingCardinality/DataHost.lean:64--65,153--154}),
\nolinkurl{TI.sE_hom_gen}, \nolinkurl{TI.trace_conj_integral},
\nolinkurl{hermitian_pairing_integral}
(\path{GL3/PHfInt.lean:79--111,145--164,203--209}), and
\nolinkurl{FL.embA_bij}, \nolinkurl{FL.sum_tau_trace}
(\path{GL3/PHfTrace.lean:34--69}).

Write $s=\sqrt d$, $\omega=(1+is)/2$; the embedded quadratic order is
$\mathbb Z+\mathbb Z\omega$ since $d$ is prime and $d\equiv3\pmod4$.
Indeed for an algebraic integer $x+yis$ with $x,y\in\mathbb Q$, trace and
norm give $2x\in\mathbb Z$ and $d(2y)^2\in\mathbb Z$; squarefreeness gives
$2y\in\mathbb Z$ and the norm modulo $4$ gives the same parity for both.
Conversely $\omega$ satisfies a monic integral quadratic.
A nonzero ideal of the quadratic order has finite index in its
rank-two integral lattice (it contains a nonzero principal ideal).
It is free of rank two; see \nolinkurl{CS.cc_basis}
(\path{GL3/PSelfDual.lean:494--543}). In such a $\mathbb Z$-basis the area of $C$ is
$A=|\delta|^2(s/2)m$, where $m$ is the absolute coordinate determinant
relative to $(1,\omega)$. It is positive since a nonzero ideal contains
a nonzero principal ideal of full rank. Every oriented determinant of two
periods is an integer multiple of $A$, and a basis attains $A$ up to sign,
as required by the infimum defining area. For $c,c'\in\mathfrak c$
write $\overline{\phi_0(c')}\phi_0(c)=(u+vis)/2$. Taking determinants
of the periods of $c'$ with those of $c$ and $c\omega$ gives respectively
$v/m\in\mathbb Z$ and $(u+v)/(2m)\in\mathbb Z$.
Thus the product divided by $m$ belongs to $\mathbb Z+\mathbb Z\omega$.
On generators of $\Lambda$ multiply this element by $\phi_0(h(S,T))$
and by $2/s$. Additivity now gives, for $x,y\in\Lambda$,
\[
 H(x,y)=a/s+ib,\qquad a,b\in\mathbb Z,\quad a\equiv b\pmod2.
\]
This proves $\Lambda\subset\Lambda^\vee$ for $\Im H$. See
\nolinkurl{AR.okc_of}, \nolinkurl{AR.div_okc}, \nolinkurl{HfArithRec_of_int}
(\path{GL3/PHfArith.lean:34--37,85--86,295--306,321--389,402--419}) and
\nolinkurl{hermitian_form_integral_values} (\path{GL3/PThIndep.lean:11}).

% blueprint/ch_P_r14.tex lines 127 through 163; r14:lat-dual
\paragraph{Recovering the dual lattice.}\phantomsection\label{r14:lat-dual}
For an oriented basis $w_1,w_2$ of $C$ the pairing
$\Im(z\bar w)/A$ has off-diagonal entries $\pm1$ on that basis.
In its real coordinates, integral pairing with the two basis vectors
therefore characterizes membership in $C$. Take $v\in\Lambda^\vee$.
Pair with all $w e(T)$, $w\in C$, $T\in M$. The preceding plane
calculation gives
$\ell(T)=\sum v_{\tau ij}\overline{\tau(T_{ij})}\in C$.
The module $M$ spans all matrices over $\mathbb Q$: full lattices $L,L'$
are commensurable with their integral coordinate lattices (clear denominators
of generators and of the expressions for the coordinate basis). Hence a
fixed nonzero integer multiplies every integral matrix into
$\operatorname{Hom}(L,L')$ and then into $M$ after one more denominator
bound for the nonzero fractional ideal $I$. Testing a single matrix entry at
$\bar y$ after clearing denominators gives
$\delta^{-1}\sum_\tau v_{\tau ij}\tau(y)\in\phi_0(K)$ for every $y\in E$.
Pulling back by $\phi_0$ defines a $K$-linear functional of $y$.
Nondegeneracy of the separable trace writes it uniquely as
$\operatorname{Tr}_{E/K}(X_{ij}y)$; independence of distinct embeddings
then gives $v_{\tau ij}=\delta\tau(X_{ij})$.
Thus $\ell(T)=\delta\phi_0(h(X,T))$ and $h(X,M)\subset\mathfrak c$.
See \nolinkurl{curve_lattice_self_dual}, \nolinkurl{SD.coord},
\nolinkurl{SelfDualRec_of_curve}
(\path{GL3/PSelfDual.lean:184--261,310--393,410--459,563--622}).

For $k\in\mathfrak c^{-1}$, $h(kX,M)\subset\mathcal O_K$.
If $x\in L$, $y\in L'$, the rank-one map $yx^*$ sends $L$ into $L'$.
Testing $a yx^*$ for every $a\in I$ gives
$\operatorname{Tr}_{E/K}(\langle y,kXx\rangle\bar a)\in\mathcal O_K$.
Trace duality of $I$ gives $\langle y,kXx\rangle\in I$.
For every $b\in\mathfrak b$ we get $bkXx\in L'$ by self-duality of $L'$,
so $bkX\in\operatorname{Hom}(L,L')$. Expressing $1$ as a sum of
products in $I\mathfrak b$ gives $kX\in M$; doing the same for
$\mathfrak c\mathfrak c^{-1}$ gives $X\in\mathfrak c M$. Consequently
$v\in\Lambda$. See \nolinkurl{SD.hom_step}, \nolinkurl{SD.rank_one_mem_hom_lattice},
\nolinkurl{SD.hYY_rank1}, \nolinkurl{SD.alg_step}, \nolinkurl{self_dual_lattice}
(\path{GL3/PSelfDual.lean:15--46,54--180,624}).

% blueprint/ch_P_r14.tex lines 165 through 187; r14:lat-indices
\paragraph{Semicharacter and the two indices.}\phantomsection\label{r14:lat-indices}
For $x\in\Lambda$, $H(x,x)$ is real, so the parity above gives
$q(x)=(s/2)\Re H(x,x)\in\mathbb Z$. Rounding fixes $q(x)$ and
$\alpha(x)=(-1)^{q(x)}$ is a sign. Polarization gives
$q(x+y)-q(x)-q(y)=s\Re H(x,y)=a$. Since $a\equiv b\pmod2$,
$\alpha(x+y)=\alpha(x)\alpha(y)\exp(\pi i\Im H(x,y))$.
For each admitted pair, $\mathfrak cM$ is commensurable with the integral
coordinate matrices: $L,L'$ have upper and lower integral denominator
bounds, hence so does the module of maps between them; multiplication by a
nonzero ideal or fractional ideal preserves such bounds by clearing the
ideal's generators and using any one nonzero member for the lower bound.
For two pairs take $R=e((\delta_0^{-1}\operatorname{Mat}(\mathcal O_E)))$
scaled by $\delta$, with a common positive integer denominator $\delta_0$,
so that $kR\subset\Lambda_1\cap\Lambda_2\subset\Lambda_i\subset R$
for some $k>0$. Then both intersection quotients are subquotients of the
finite group $R/kR$; both indices are nonzero, also in rank zero.
This proves the two clauses of \nolinkurl{Base.LatAll}
(\path{KernelLattices/LateralTheory.lean:185--189,207--213}).
See \nolinkurl{IO.semiChar_of_HfInt}, \nolinkurl{IO.La_index_finite},
\nolinkurl{IO.La_index_ne_zero} (\path{GL3/PIopIso.lean:208--267,275--293,315--338}),
\nolinkurl{FL.YY_lower}, \nolinkurl{FL.La_upper}, \nolinkurl{FL.La_lower}
(\path{GL3/PLinkThU.lean:515--541,547--608,609--634}) and
\nolinkurl{Base.latAll} (\path{KernelLocalFactors/IsAdmissibleFlagPairOfUnitriangularLeftAutomorph.lean:623--644}).

% blueprint/ch_P_r14.tex lines 189 through 270; r14:line
\paragraph{Why comparison on a theta line gives a scalar.}\label{r14:line}
At the admitted pair let $A$ be the period lattice and let $H$ be the
area-normalized Hermitian form, linear in its first argument. The reference
and natural vectors have the same automorphy factor on $A$. We need
dimension at most one before a nonzero reference vector can serve as a
basis. For the Hom lattice, take arithmetic trace-dual families $y_i,y_i'$
with $x=\sum_i h(x,y_i')y_i$. After applying the embeddings this expansion
also holds over $\mathbb C$: distinct embeddings are linearly independent
and their evaluation vectors span each matrix coordinate. The coordinates
of any element of $A$ against $y_i'$ belong to the planar period lattice.
They determine the element by the expansion, so their map embeds $A$
continuously into a finite power of a discrete lattice and $A$ is
discrete. Its imaginary dual equals $A$ by the lattice argument above.
For the alternative orthonormal presentation of the lattice the same
expansion holds with $y_i'=y_i$.
See \nolinkurl{ThetaLine.expand}, \nolinkurl{ThetaLine.of_dual_families}
and \nolinkurl{ThetaLine.theta_line}
(\path{FoundationProofs/ThetaSpaceSelfDualPairAtMostLine.lean:82--132,145--169,187--219,276--294}).

Identify the coordinate space with $\mathbb C^g$, where
$g=|\mathrm{Tau}|NN'$. Multiplying by the inverse of the positive
area gives the positive definite Hermitian matrix. For $a\in A$ an entire
theta function satisfies
\[
 f(z+a)=\alpha(a)^{-1}\exp\big(\pi H(z,a)+\tfrac\pi2 H(a,a)\big) f(z).
\]
The sign from interchanging the arguments of $\operatorname{Im}H$ preserves
the integral-dual condition.
The integral alternating form on $A$ is perfect: a real linear functional
is represented by pairing with some $y$ by nondegeneracy, and if it takes
integral values on $A$, self-duality puts $y$ in $A$. Uniqueness follows
from real spanning and positive definiteness. Indeed a nonzero vector
symplectically orthogonal to the real span would put an entire real
line in the discrete integral dual. Thus the lattice has real span
$\mathbb C^g$ and integral rank $2g$. The same argument works for
the frame lattice: there the orthonormal frame is its own dual family
in \nolinkurl{ThetaLine.of_dual_families}.
A primitive basis vector in the perfect alternating lattice pairs to
one with a second vector. Split off their unimodular plane and iterate
on its orthogonal complement to obtain an integral symplectic basis.
See \nolinkurl{ThetaTranslSpan.span_eq_top}
(\path{UnitaryExistence/ExistsImEqLattice.lean:106--122}) and
\nolinkurl{SiegelNormalForm.exists_normal_form}
(\path{ToricGeometry/IsUnitOfLagrangian.lean:177--353}); its linear
exponential term uses logarithms of the unit phases on the first
basis vectors. In a symplectic integral
basis of this principal form, the first half gives a complex basis by
positivity on a Lagrangian. In the resulting coordinates, multiplication
by the nonvanishing exponential of a quadratic polynomial makes the
first $g$ periods ordinary periods. The second half then acts on the
transformed function $u$ by
$u(z+\tau_j)=b_j\exp(-2\pi i z_j)u(z)$ with $b_j\ne0$.
These are the degree-one transformation laws obtained in
\nolinkurl{ThetaCoordLine.transl_of_mem}, \nolinkurl{ThetaCoordLine.exists_eq_smul}
(\path{FoundationProofs/ThetaFunctionSelfDualExistsEqSmul.lean:25--52,54--93}),
\nolinkurl{ThetaTranslSpan.exists_eq_smul}
(\path{UnitaryExistence/ExistsImEqLattice.lean:162--235}), and
\nolinkurl{card_le_pow_of_linearIndependent_theta_principal}
(\path{ToricGeometry/CardLePowOfLinearIndependentThetaPrincipal.lean:11--59}).

Let $c_m(u)$ be the Fourier coefficient on the real unit cube, for
$m\in\mathbb Z^g$. Translating its contour by $\tau_j$ is valid coordinate
by coordinate by holomorphy and cancellation of the opposite faces.
Equating the two evaluations gives
\[
 b_jc_{m+e_j}(u)=\exp(2\pi i\,m\cdot\tau_j)c_m(u).
\]
Both factors are nonzero. Thus $c_0(u)=0$ forces every coefficient to vanish.
For every complex translation of the real unit cube the translated
coefficients are the original ones multiplied by nonzero exponentials.
Fourier uniqueness on each translated real torus gives zero there, hence
zero at every point of the complex space.
This works for $g=0$ as well, when evaluation is the only coefficient.
The map to $c_0$ is injective on these sections. For a nonzero reference
vector $G$, subtracting the right multiple of $G$ from $F$ now proves
$F=aG$ before undoing the invertible change of coordinates.
See \nolinkurl{fourierCoeff_unitCube_translate_of_entire_periodic}
(\path{ToricGeometry/FourierCoeffUnitCubeTranslateOfEntirePeriodic.lean:70--135}),
\nolinkurl{entire_periodic_eq_zero_of_fourierCoeff_eq_zero}
(\path{ToricGeometry/EntirePeriodicEqZeroOfFourierCoeffEqZero.lean:26--96})
and \nolinkurl{ClassicalTheta.card_le_pow_of_linearIndependent}
(\path{ToricGeometry/CardLePowOfLinearIndependent.lean:50--120}).

% blueprint/ch_P_r14.tex lines 272 through 576; r14:vec, r14:vec-transport, r14:vec-fixed, r14:vec-descent
\subsection{The natural nonzero vector and its field of coefficients}
\label{r14:vec}
Fix \(D,\mathfrak c,\delta,B,\chi\) from the tied entry, put
\(T=B.\mathrm{tor}=\mathrm{tor}_0(D,\mathfrak c,\delta)\), and write
\(\Lambda=\delta\varphi_0(\mathfrak c)\).
The entry gives a curve ideal, a twelfth root of its lattice discriminant,
and a character multiplicative on norm-one ideals with prescribed values on
principal unitary ideals. For each pair of even ranks \(N,N'\) and
self-dual lattices \(L,L'\) we will exhibit, for the prescribed function
\(F=B.\mathrm{eNat}(\chi)_{L,L'}\),
\[
 F\in\mathrm{Th}(D,T,L,L'),\qquad F\ne0,\qquad
 \exists M\subset\mathbf C,\quad [M:\mathbf Q]<\infty,\quad
 \forall\alpha,\quad a_\alpha(F)\in M .
\]
Here \(a_\alpha\) is the derivative at zero in the indicated coordinate
directions divided by \(\prod_k\alpha(k)!\). The coordinates are
\((\tau,i,j)\) with \(i\in\mathrm{Fin}(N')\), \(j\in\mathrm{Fin}(N)\).
The order of the two final quantifiers permits \(M\) to depend on the
pair. See \nolinkurl{TaylorLogMetric.vec22} and
\nolinkurl{Base.VecP} (\path{GlobalSchurBounds/TaylorLogMetric.lean:13--16},
\path{KernelLattices/LateralTheory.lean:215}).

\paragraph{Construction and nonvanishing.}\label{r14:vec-transport}
Write \(S,S'\) for the standard lattices of the indicated ranks,
\(A_0=\mathrm{La}(S,S')\), \(A_1=\mathrm{La}(L,S')\) and
\(A_2=\mathrm{La}(L,L')\). Write \(I_{ij}\) for the normalized
intertwiner from \(A_i\) to \(A_j\). On the target pair put
\[
 \begin{aligned}
  C_\chi&=\chi(\operatorname{detId}(S,L))^{N'/2}\,
           \chi(\operatorname{detId}(S',L'))^{-N/2},\\
  V_0(z)&=\prod_{i,j}\theta_M\big((z_{\tau i j})_\tau\big),
     & F&=C_\chi I_{12}(I_{01}V_0).
 \end{aligned}
\]
The symbol \(\operatorname{detId}\) denotes the ideal spanned by the determinants
of matrices carrying the first lattice into the second. The vector
\(\theta_M=B.\mathrm{thMM}\) is the chosen theta function whose
fourfold tensor product is \(B.\kappa\) times the standard fourfold
kernel in the coordinates \(B.g_C\)
(\path{KernelLattices/LateralTheory.lean:23--36,145--179}).
The supplied period pair has lattice \(\Lambda\) and positive area.
Its half-period and the two semicharacter comparisons, self-duality,
and both finite indices give \(B.\mathrm{StdOK}\)
through the inputs \nolinkurl{SO.stdOK} and \nolinkurl{FourthRootStd.stdOK}
(3,161 and 2,853 inclusive theorem lines beyond the common host baseline,
respectively; the same convention applies to supplier sizes below).
See \path{GL3/PStdOK.lean:326--339} and \path{SpanRangeEmb4EqTop.lean:700--701}.
Each local root in \(B.\kappa\) has a nonzero prescribed power; its
finite product with integral exponents is nonzero. The fourfold identity and nonvanishing
use \nolinkurl{thM_spec} and \nolinkurl{isThM_ne_zero} with $B.\mathrm{StdOK}$
and $B.\kappa\ne0$ (3,001 lines;
\path{KernelLocalFactors/IsAdmissibleFlagPairOfUnitriangularLeftAutomorph.lean:75--81,218--220}).
The earlier scalar calculation is in Paragraph~\ref{par:r4-z43-iso};
Z4.3 is partial, and review A treated this as a supplier. At a point with all columns
equal to a point of nonvanishing, \(V_0\ne0\); the empty product
covers zero rank. The tensor-product translation rule on \(A_0\)
uses the sum of the integral quadratic terms of its columns, giving
\(V_0\in\mathrm{Th}(A_0)\).
The identification with the tensor product lattice and the equality of
the two area normalizations enter via \nolinkurl{Base.lat0OK},
\nolinkurl{Base.areaOK} and \nolinkurl{Base.thZeroMem}
(\path{KernelLocalFactors/IsAdmissibleFlagPairOfUnitriangularLeftAutomorph.lean:318--319,438--439,480--501}).

The three pairs are admitted. For each of the two adjacent pairs,
\(\mathrm{LatAll}\) supplies self-duality and the same semicharacter
on both lattices and both finite relative indices. Averaging the target
signed translations over representatives sends a theta vector to a
theta vector; on the source fixed space the reverse normalized average
composed with the forward one is the identity. Indeed, summing a
nontrivial character of the finite intersection quotient gives zero,
and its trivial character gives its order, which cancels the two
normalizations: the alternating form pairs the two intersection
quotients perfectly into $\mathbf Q/\mathbf Z$, so their orders coincide.
Indeed their annihilators are the respective intersections by self-duality.
Under $B.\mathrm{StdOK},B.\mathrm{LatAll},B.\mathrm{ThZeroMem}(N,N')$
and admission, both transports preserve nonvanishing.
The determinant ideals are norm one by
\nolinkurl{relative_determinant_norm_one}
(\path{DeterminantNormOneOfUnitDisc.lean:409}). For each such ideal a positive
power is generated by a unitary element, so the corresponding power
of its character is a product of inverses of CM conjugates of that
element. All are algebraic of absolute value one. Both factors of
\(C_\chi\) are therefore nonzero and algebraic for the displayed
integral exponents. Scalar multiplication preserves the theta relations
and yields the first two conclusions. The source implements these steps
in \nolinkurl{Base.thMoved_ne_zero}, \nolinkurl{Base.vec_mem_ne}
(\path{KernelLocalFactors/IsAdmissibleFlagPairOfUnitriangularLeftAutomorph.lean:197--222,234--271,480--508,646--660}),
and \nolinkurl{GlobalTheory.MomT.FromComb.chiFac_mem}
(\path{GlobalInterpolation/TranslateMomentsAllOfJetMomR.lean:323--339}).

\paragraph{The line and the reference vector.}\label{r14:vec-reference}
The CM embeddings restrict to \(\varphi_0\) on \(K\), and their sum
on \(E\) is \(\varphi_0\circ\mathrm{Tr}_{E/K}\): the indexing set
enumerates \(\mathrm{Gal}(E/K)\) once. This is the
\nolinkurl{admG1} used in \nolinkurl{GlobalTheory.ThetaLine.theta_line}.
The trace-dual coordinate argument and the determining Fourier coefficient
in Section~\ref{r14:line} show that any two such vectors are proportional
once the reference vector is nonzero.
See \nolinkurl{GlobalTheory.InducedType.induced_type}
(\path{FoundationProofs/InducedTypeThetaGenusOnePermutes.lean:43--53,74--75})
and \nolinkurl{GlobalTheory.ThetaLine.theta_line}
(\path{FoundationProofs/ThetaSpaceSelfDualPairAtMostLine.lean:276--294}).

For the same curve put \(n=\mathrm N\mathfrak c\), which is odd, and
\(t_0=\delta n/2\). Then \(2t_0\in\Lambda\) since \(n\in\mathfrak c\).
If either \((1\pm i\sqrt d)t_0/2\) belonged to \(\Lambda\), dividing
by \(\delta\ne0\) would give an algebraic integer of trace \(n/2\)
over \(\mathbf Q\), impossible for odd \(n\). Write
\(A=\mathrm{area}(\Lambda)\) and let \(\sigma^*_\Lambda\) be the modified
sigma function \nolinkurl{sgStar}, normalized by \((\sigma^*_\Lambda)'(0)=1\).
The genus-one vector uses the exponential Heisenberg translation
\[
 \begin{aligned}
 (U_1(t)h)(z)
    &=h(z+t)\exp\!\left(-\frac{\pi z\overline t}{A}
                          -\frac{\pi t\overline t}{2A}\right),\\
 \mathrm{th}_0(z)
    &=\frac{(U_1(t_0)\sigma^*_\Lambda)(z)}
            {(U_1(t_0)\sigma^*_\Lambda)(0)} .
 \end{aligned}
\]
Here \(t_0\notin\Lambda\) by the unkilled half-period condition and
the stability of \(\Lambda\) under the CM integers. Hence the modified
sigma function is nonzero there: \nolinkurl{OddTheta.sgStar_laws}
constructs the normalized entire function whose zeros are exactly the
period lattice (\path{OddThetaFunctionExistenceUniqueness.lean:115--147}).
Consequently \(\mathrm{th}_0(0)=1\)
(\path{GlobalSchurBounds/GL3Theta.lean:196--203},
\path{GeneralLinearActions/IsDiscriminantTwelfthRootOfIdealLattice.lean:234--241}).
For even \(NN'\) choose an orthonormal
frame for the standard pair at 2. Pull back its product of genus-one
thetas and apply the normalized intertwiner from the frame lattice
to \(A_0\), with the algebraic scalar of absolute value one prescribed
by the vacuum construction. The product at the origin is one.
Applying the reverse intertwiner detects its nonvanishing. The resulting
\(e_0\) satisfies the standard vacuum conditions \nolinkurl{oddE0};
these include its description by that frame and scalar. The construction
uses \nolinkurl{GlobalTheory.TwoTorsionPt.exists_point} and
\nolinkurl{GlobalTheory.Vacuum.exists_vacuum}
(\path{FoundationProofs/ExistsTwoTorsionCMPointAllConjugates.lean:16--99,127--130},
\path{FoundationProofs/NormSqMulGaussianLatticeTranslateEqShifted.lean:276--355}).
The latter is an input of 16,557 lines for the full vector family.
In its frame at 2, \nolinkurl{OddGenus.odd_one} first obtains a rational
orthonormal frame for a positive diagonal Hermitian form with square
determinant (\path{StandardPairLatticeLocallyOrthonormalBasis.lean:138,337--357,420--425}).
The quadratic local-to-global input \nolinkurl{Published.HasseNormSt}
takes $d,x\ne0$ with $x=a^2-db^2$ at every finite and real place and
produces such a representation over the base number field. In the split
case $d=r^2$, use $a=(x+1)/2$, $b=(x-1)/(2r)$ as in
\nolinkurl{Published.HasseNormSt.of_isSquare}. For nonsquare $d$,
\nolinkurl{QNorm.hasseNormSt} supplies the local-to-global norm step
(\path{ExistsNormTensor.lean:12--25},
\path{HasseNormTheoremQuadraticExtension.lean:13--32}).
The comparison with Hasse's cyclic norm theorem, including the
complex places, is on p.~\pageref{r15:hasse-comparison}.
In particular \nolinkurl{SqDet.exists_rat_sum_norm_forms_eq_one}
uses isotropy of the rational form $\langle a,ad,b,bd,-1\rangle$
to produce a vector of norm one; see
\path{ExistsOrthonormalHermitianOfDiagPosIsSquare.lean:17--54},
\nolinkurl{LocalQuad.isotropic_of_indefinite}
(\path{ExistsCongrIffTraceFormsIsometricCm.lean:290--298}) and
\nolinkurl{LocalQuad.exists_norm_of_quaternary}
(\path{ExistsQuaternaryOfCoordinates.lean:14--40}).
These are the three literature-boundary declarations on this frame route.

In \nolinkurl{GlobalTheory.TrueTheta.exists_family_of_curve} a character
\(\chi'\) descends from an idelic character \(\kappa'\) trivial on
integral units, with its value on a relative determinant ideal read
from an adelic mover. Choose movers \(g,g'\) from the standard
lattices to \(L,L'\). Their unit scalar \(s\), given by those two
determinant factors with exponents \(N'/2,-N/2\), is algebraic by
\nolinkurl{HeckeVal.hecke_values}
(\path{UnitaryHeckeCharacterValues.lean:232--278}) in the vacuum scalar
comparison. Thus the reference vector for this pair is
\(e=sG\), \(G=I_{02}e_0\), with a \emph{single} direct transport.
It is a nonzero theta vector by the inverse-intertwiner identity.
Its use here requires only \(s\in\overline{\mathbf Q}^{\times}\);
one could take \(s=1\) in the coefficient-field argument for this
individual pair. The reference character and the natural character
can differ. See \path{IwasawaDecomposition/ExistsThetaVectorFamilyOfCurveDatum.lean:15--52,75--90}
and \path{IwasawaDecomposition/ExistsMovesOfSelfDual.lean:77--135,570--581,659--675}.

\paragraph{A number field for the reference vector.}\label{r14:vec-fixed}
Here is the finite descent inside
\nolinkurl{GlobalTheory.TrueTheta.exists_coef_field}. The CM period
calculation gives algebraic \(g_2(\Lambda),g_3(\Lambda)\): the normalized
discriminant and the integral CM \(j\)-invariant first determine
\(g_2^3\), then \(g_3^2\). The resulting CM theta calculation gives
algebraic Taylor coefficients of \(G\). More explicitly, the
genus-one normalized theta and its rational translates have algebraic
jets; products in the orthonormal CM coordinates preserve this property.
A normalized intertwiner is a finite sum of rational translates with
algebraic signs and inverse square root of an integer. The Heisenberg
phase arising on composing rational translates is a root of unity.
Hence the same algebraicity holds after each intertwiner. The needed
source laws are \nolinkurl{CurveAlgebraic.integral_j_and_algebraic_data},
\nolinkurl{ThetaAlgebraic.taylor_algebraic}, and the conjugation law
\nolinkurl{ThetaConj.useConj_all}. For every algebraic embedding of
the coefficient field the last law takes a transported vacuum to a
signed transport of the conjugate vacuum, on the conjugate torus.
The sign is nonzero. These laws use the curve and its normalized
half-period, the vacuum's frame description, self-duality of the
input and output pairs, and even \(NN'\)
(\path{FoundationProofs/IntegralJAndAlgebraicData.lean:126--151},
\path{FoundationProofs/ThetaTaylorCoeffsAlgebraicAtComplexMultiplicatio.lean:594--616},
\path{FoundationProofs/ThetaVacuaConjugatesCmPoints.lean:529--567}).

Choose \(\alpha_0\) with \(a_{\alpha_0}(G)\ne0\) by the Taylor expansion
of the nonzero entire \(G\). For a rational basis \(b_j\) of \(E\),
adjoin to \(\mathbf Q\) the finite set
\[
 \{\mathrm{emb}_\tau(b_j)\}_{\tau,j}
 \ \cup\ \{g_2(\Lambda),g_3(\Lambda),s,a_{\alpha_0}(G)\}.
\]
Call the resulting number field \(M_0\). Each automorphism \(\sigma\) of
\(\overline{\mathbf Q}/M_0\) fixes every CM embedding on that basis
and fixes \(g_2(\Lambda),g_3(\Lambda)\). The torus constructed from the
conjugated embeddings and the lattice with those invariants is
therefore \(T\) by uniqueness of the lattice with fixed $(g_2,g_3)$,
\nolinkurl{ThetaConj.lamConj_eq_self}
(\path{CmThetaTaylorCoeffAlgebraicAndConjugate.lean:47--59}).
Coefficientwise conjugation of \(G\) gives an
entire theta vector on \(A_2\) by the conjugation law. It is nonzero:
reverse the signed transport; the resulting conjugate vacuum still
has the normalized half-period product description. The line argument
makes this vector \(aG\). Its holomorphy and nonvanishing justify recovering
the coefficients of its defining conjugated series by
\nolinkurl{ConjSeries.taylor_fnConj}
(\path{TaylorCoeffConjugateSeries.lean:207--223}).
Coefficients of the conjugate are
\(\sigma(a_\beta(G))\), and linearity gives \(a\,a_\beta(G)\)
for every \(\beta\). Comparing at \(\alpha_0\), whose nonzero
coefficient is fixed, gives \(a=1\). All \(a_\beta(G)\) belong to
the fixed field \(M_0\), by the infinite Galois correspondence
\nolinkurl{InfiniteGalois.fixedField_fixingSubgroup} applied to the finite
set of generators (\path{ExistsMovesOfSelfDual.lean:407--420}), hence so do all \(a_\beta(e)\). This uses
\nolinkurl{GlobalTheory.TrueTheta.exists_field} and
\nolinkurl{GlobalTheory.TrueTheta.exists_coef_field}
(\path{IwasawaDecomposition/ExistsMovesOfSelfDual.lean:407--440,442--515}).

\paragraph{The algebraic scalar for the prescribed natural vector.}
\label{r14:vec-descent}
We also need algebraicity on the natural side. The standard
fourfold tensor identity expresses the fourfold product of
\(\theta_M\) as an algebraic multiple of a finite sum of products
of rationally translated genus-one functions in algebraic linear
coordinates. Choose an index $\beta$ of least total degree with nonzero
normalized Taylor coefficient $a_\beta$ of $\theta_M$; one exists by
entireness and nonvanishing. The four independent blocks have coefficient
$a_\beta^4$ at $(\beta,\beta,\beta,\beta)$ and $a_\alpha a_\beta^3$
at $(\alpha,\beta,\beta,\beta)$ for every $\alpha$. Each is algebraic by the
fourfold identity, hence $a_\beta$ is algebraic and nonzero; division
gives $a_\alpha\in\overline{\mathbf Q}$ for all $\alpha$. Products of
these expansions give the same conclusion for $V_0$.
The input \nolinkurl{AlgTranslate.coefIn_translate}
(\path{CoefInTranslate.lean:53--85}) applied in the standard theta
calculation gives
algebraic jets of \(U_x V_0\) for every \(x\) in the rational saturation
of \(A_0\). Every period of \(A_1,A_2\) belongs to this saturation.
Under successive translates the phase is the exponential of \(\pi i\)
times a rational number. Therefore the two finite signed averages
preserve this translated-jet property, and so does multiplication by
\(C_\chi\in\overline{\mathbf Q}\). This last deduction uses the supplier
\nolinkurl{GlobalTheory.ThLoc.algTAll} (15,462 lines), which gives algebraic Taylor
coefficients of \(U_xF\) for every \(x\) with \(m x\in A_2\) for
some positive integer \(m\); set \(x=0,m=1\), for which \(U_0F=F\).
See \path{IwasawaDecomposition/VacuumEntireAndNormalizedJetsAlgebraic.lean:121--149,197--258,280--318}
and \path{GaloisGeometry/AlgTAll.lean:12--27}.

Both \(F\) and \(e\ne0\) lie on the same theta line on \(A_2\);
write \(F=be\). Then \(a_\beta(F)=b\,a_\beta(e)\) for every
multi-index by complex linearity of the derivative. Choose \(\alpha\)
with \(a_\alpha(e)\ne0\). Its numerator \(a_\alpha(F)\) is algebraic
and the denominator belongs to \(M_0\), so
\(b=a_\alpha(F)/a_\alpha(e)\in\overline{\mathbf Q}\).
The compositum \(M=M_0\mathbf Q(b)\) is finite. It contains all
\(a_\beta(F)\) simultaneously. Moreover \(b\ne0\) since \(F\ne0\).
This scalar is chosen anew for each pair. The comparison for all
pairs of fixed ranks in the equivariance proof requires a separate
argument (\nolinkurl{GlobalTheory.ColNorms.familyVectors},
\path{DivisorProofs/NormalisedFamilyNonzeroThetaVectorsNumberField.lean:15--50}).

\paragraph{Interface at the tied entry.}
The formal family theorem additionally takes the host field conditions,
the induced CM type, and finite support of the relative determinant
exponents of \(B.\mathrm{kc}\). It eliminates these three explicit
arguments before starting its proof at a pair. In the call at the
tied entry they have independent suppliers. The host field conditions
come from the host ramification theorem, and the restriction of each
CM embedding to \(K\) holds because its indexing automorphism fixes
\(K\). The ignored finite-support input has its own proof.\footnote{For
$\mathrm{yFinite}$ clear denominators of the fixed map and its inverse
on finite integral generators. At any split prime beyond the product
of the two denominators both maps are integral at the selected prime,
so the relative determinant is a unit.} These are
\nolinkurl{GL3Theta.cm_field_facts_hold},
\nolinkurl{GL3Theta.induced_cm_type}, and
\nolinkurl{GlobalTheory.ExponentsFinite.yFinite}
(\path{UnitTraceSumBounds/InducedCmType.lean:10--17},
\path{DivisorProofs/KernelConstantExponentsFinite.lean:85--126}).
The proof supplies the prescribed \(B.\mathrm{VecP}(\chi)\) using the
same base and character as the other projections of
\nolinkurl{GlobalTheory.TieBase.base_data} conditional on the indicated suppliers.

% blueprint/ch_P_r14.tex lines 578 through 1058; r14:equivariance, r14:eq-reference, r14:eq-transport, r14:eq-rational, r14:eq-normalize, r14:weil-sign
\subsection{Equivariance of the normalized theta family}\label{r14:equivariance}
Fix a curve ideal $\mathfrak c$, a discriminant twelfth root $\delta$ and a
base $B$ on these data. Write $T=B.\mathrm{tor}$, and take any $\chi$
satisfying $\mathrm{IsChi}(D,\chi)$. For even ranks $n,m$ and self-dual
Hermitian lattices $L,M$ of these ranks, put $A_{LM}=\mathrm{La}(D,T,L,M)$,
$S=\mathrm{stdL}(D,n)$ and $S'=\mathrm{stdL}(D,m)$. Write $I_{A,B'}$ for the
normalized translation average and $e_\chi(L,M)=B.eNat(\chi,n,m,L,M)$.
For $g,h$ rational unitary matrices define
$P_{g,h}F(z)=F((h_\tau^*z_\tau g_\tau)_\tau)$, using the torus's
embeddings. The three assertions are
\begin{align*}
 P_{g,h}e_\chi(L,M)&=e_\chi(gL,hM),\\
 e_\chi(L,M_1)&=\chi(\operatorname{detId}(M,M_1))^{-n/2}\,
 I_{A_{LM},A_{LM_1}}e_\chi(L,M),\\
 e_\chi(L_1,M)&=\chi(\operatorname{detId}(L,L_1))^{m/2}\,
 I_{A_{LM},A_{L_1M}}e_\chi(L,M).
\end{align*}
Here and below $\operatorname{detId}$ denotes the relative determinant ideal. The exponent
in the second line uses the unchanged source rank, and the one in the third
uses the unchanged target rank. Zero ranks are allowed. See
\nolinkurl{Base.EqP} (\path{KernelLattices/LateralTheory.lean:219}),
\nolinkurl{eq22} (\path{GlobalSchurBounds/TaylorLogMetric.lean:37--50}) and
\nolinkurl{pull} (\path{IsDiscriminantTwelfthRootOfIdealLattice.lean:106--108}).

\paragraph{One reference character and one family.}\phantomsection\label{r14:eq-reference}
For a finite idele $u$ write $J(u)=\{a\in E:u^{-1}a\in\widehat{\mathcal O}_E\}$.
Two ideles representing the same fractional ideal differ by an integral
unit. The character $\kappa$ in \nolinkurl{TrueTheta.exists_kappa} is trivial
on these units. Hence $\eta(J(u))=\kappa(u)$ is independent of $u$.
Set $\eta$ equal to zero on unrepresented submodules. Every norm-one
ideal has a representative, $J(uv)=J(u)J(v)$, and a rational
norm-one element $b$ represents $(b)$ with value
$\prod_\tau\tau(b)^{-1}$. These give both clauses of $\mathrm{IsChi}$.
If $g_1,g_2$ move the standard lattice to $L_1,L_2$, the criterion for
membership in $\operatorname{detId}(L_1,L_2)$ is integrality after division by
$\det(g_1^{-1}g_2)$. Thus
$\eta(\operatorname{detId}(L_1,L_2))=\kappa(\det(g_1^{-1}g_2))$.
The particular $\kappa$ is shared across the ranks. Its construction
restricts a Hecke character to the split-place idempotent $e$ by
$u\mapsto eu+(1-e)$. The splitting characters have half-weight parameters $k=-n/2$ and
$k'=m/2$ and infinity types $2k=-n$ and $2k'=m$, respectively.
They give a Weil operator $W$ with
$\epsilon n_\tau(\psi)+n=0$ and
$\epsilon' n_\tau(\psi')-m=0$. On the split factor its vacuum multiplier
is $\psi(y')^{\epsilon}\psi'(y)^{\epsilon'}$, with
$y/\bar y=\det(g_{\rm split})$ and
$y'/\bar y'=\det(h_{\rm split})$. The inert sign in the product formula
has even exponent. The resulting multiplier, for $g,h$ moving $S,S'$ to
$L,M$, is
\[
 c(g,h)=\kappa(\det g)^{m/2}\kappa(\det h)^{-n/2}
       =:C_\eta(L,M).
\]
Three inputs in \nolinkurl{TrueTheta.exists_kappa} are needed before the
vacuum identity. \nolinkurl{TrueTheta.exists_base} extends a unitary
character trivial on integral ideles and supplies Hecke characters with
weights $2k$ (402 lines, \path{ExistsMovesOfSelfDual.lean:519--564}).
\nolinkurl{InertTrivial.character_trivial_at_inert} removes the inert
part on the rational norm-one idele; this gives the value
$\prod_\tau\tau(b)^{-1}$ required for the reference character in
\nolinkurl{NatEqui.chi_smulLat} (6,020 lines,
\path{CharacterTrivialAtInertPrimes.lean:255--309},
\path{ExistsMovesOfSelfDual.lean:592--607}). \nolinkurl{HeckeVal.hecke_values}
gives algebraicity and conjugate modulus one of the scalar factors
(879 lines, \path{UnitaryHeckeCharacterValues.lean:232--278}).
The operator inputs are \nolinkurl{FinWeil.exists_weil},
\nolinkurl{InertSign.vacuum_inert} and \nolinkurl{VacuumExact.vacuum_exact}
(\path{ExistsWeilOfVacuum.lean:82--83},
\path{WeilRepresentationVacuumSignInertPlaces.lean:254--327},
\path{WeilOperatorVacuumExactConstantEvenRank.lean:79--105}).
Review A of round 14 reported 5,355 theorem lines in its ordered partition. The order-independent three-root cut in Section~\ref{r15:supplier-classification} gives 5,132. See \nolinkurl{TrueTheta.idele_ideal_laws},
\nolinkurl{mem_detId_iff}, \nolinkurl{scal}, \nolinkurl{exists_kappa} and
\nolinkurl{exists_chi}
(\path{IwasawaDecomposition/ExistsMovesOfSelfDual.lean:137--230,404--406,542--556,566--657,659--695}).

Choose the distinguished two-torsion point and a nonzero standard vacuum satisfying \nolinkurl{oddE0}
$v$ for each even rank product. Define
$E(L,M)=C_\eta(L,M)I_{A_{SS'},A_{LM}}v$. Since the character value is
defined on the ideal, arbitrary choices of adelic movers leave $E$ fixed.
Self-dual lattices admit such movers. In the source, a common denominator
$m_0$ bounds a generating set of $L$. A rational unitary change puts
a multiple $sL$ in standard coordinates, with $s$ coprime to $m_0$.
This is the input \nolinkurl{RepEqAt.repExS} (5,038 lines,
\path{ExistsTriv.lean:54--91}). Five boundary declarations in its
Hermitian genus argument have the following uses. Column pairings are
the entries of $P^tG\bar P$ by \nolinkurl{Published.herm_col};
\nolinkurl{Published.cc_ne_zero} and \nolinkurl{Published.det_map_cc}
give $\overline{x}\ne0$ for $x\ne0$ and
$\det(\bar P)=\overline{\det P}$. For the quadratic space over the
maximal real subfield, \nolinkurl{Published.exists_cc_eq_neg} gives
$\theta\ne0$ with $\bar\theta=-\theta$, and
\nolinkurl{Published.exists_basis_one_imaginary} gives the basis
$1,\theta$. See \path{AlgebraMapNormOfBasisOneSqrt.lean:33--36,50--57,252--277}.
These enter \nolinkurl{SelfDualGenus.exists_coord},
\nolinkurl{SelfDualGenus.exists_of_locConj_eq} and the Gram-matrix and
orthogonal-frame steps in \path{FiniteResidueFieldAdicCompletionIntegers.lean:109--129,376--446,639--703}.
Their role as elementary inputs to the genus argument is recorded on
p.~\pageref{r15:hasse-comparison}.
Using complementary conjugation-invariant idempotents divisible by $s$
and $m_0$ respectively, patch this unitary with the identity. The two
integrality conditions put $L$ inside the resulting moved lattice.
Both are self-dual, so dualizing reverses this inclusion and gives
equality. The character and vacuum normalization give
$W(g,h)v=E(L,M)$ for these movers. The average preserves theta membership
and its reverse recovers $v$, so $E$ is nonzero. Its coefficient number field
uses the algebraicity and conjugation of the averaged jets. Adjoin the
images under every $\tau$ of a rational basis of the CM field, the two
invariants $g_2$ and $g_3$ of the elliptic lattice,
$c(g,h)$, and one nonzero Taylor coefficient of the average. An automorphism
fixing these data preserves its theta line by conjugation, and comparing
the chosen coefficient forces its proportionality scalar to be one.
Every Taylor coefficient, already algebraic, is fixed and belongs to the
same finite extension of $\mathbf Q$. See \nolinkurl{TrueTheta.exists_moves_of_selfDual},
\nolinkurl{exists_coef_field}
(\path{ExistsMovesOfSelfDual.lean:77--135,407--420,442--515}) and
\nolinkurl{TrueTheta.exists_family_of_curve}
(\path{ExistsThetaVectorFamilyOfCurveDatum.lean:15--46,68--81}).

\paragraph{Transport between two moved pairs.}\phantomsection\label{r14:eq-transport}
The full intertwining statement takes a curve datum with $\delta\ne0$,
ranks $n,m$, a standard vector and a Weil operator $W$ for it.
Use paired notation $g=(g,g')$, $h=(h,h')$, $k=(k,k')$.
Suppose the two components of each of $h,gh,k,gk$ move their respective
standard lattices to the components of $L_1,L_2,M_1,M_2$ (eight move
hypotheses). For every $F\in\mathrm{Th}(L_1)$,
\[
 W(g) I_{A(L_1),A(M_1)}F
   = I_{A(L_2),A(M_2)} W(g)F.
\]
Here $A(P)$ is the period lattice of the pair $P$. In
\nolinkurl{WeilInter.om_iop} the three inputs are the coset correspondence
\nolinkurl{exists_thetaLattice_correspondence_weilOperator} (1,341 lines),
membership after the move \nolinkurl{WeilTheta.om_mem_Th_moved} (5,583),
and all-level closure \nolinkurl{ThetaLevels.levels_closed} (4,010).
See \path{ExistsThetaLatticeCorrespondenceWeilOperatorStatement.lean:7--8},
\path{WeilRepresentationThetaMovedLattice.lean:21--43},
\path{ThetaFunctionsAllLevelsClosure.lean:29--68} and
\path{WeilRepresentationCommHeisenbergIntertwiner.lean:88--100}.
Let $(g,h)$ and $(g_1,h_1)$ move the standard pair to $(L,M)$ and
$(L_1,M_1)$. Apply the vacuum identity first at the relative movers
$k=g^{-1}g_1$, $k'=h^{-1}h_1$. Acting next by $W(g,h)$ pushes their
average to the average from $A_{LM}$ to $A_{L_1M_1}$. Indeed the lattice
correspondence identifies the two finite quotients of the target lattice
by its intersection with the source lattice. It therefore preserves the
factor $[A':A'\cap A]^{-1/2}$. On each coset the Weil action carries a
Heisenberg translate to the corresponding translate; changing its
representative by an element of the intersection preserves the
semicharacter-weighted summand. Summation proves the intertwining identity.
Related inputs are the lattice facts used for coordinate moves
(\nolinkurl{VacuumScalar.lattice_facts}, 4,615 lines,
\path{OddnessConditionHolds.lean:378--414}), unitarity of the normalized
intertwiner (\nolinkurl{theta_translation_intertwiner_unitary},
3,407 lines, \path{ThetaTranslationIntertwinerUnitary.lean:183--202}),
and existence of an arithmetic lattice for each adelic mover
(\nolinkurl{NormSum.exists_lattice_of_adelic_point}, 1,105 lines,
\path{ExistsLatticeOfAdelicPoint.lean:231--252}).
The representation product law now gives
\[
 E(L_1,M_1)=c(k,k')I_{A_{LM},A_{L_1M_1}}E(L,M).
\]
For $g_1=g$ the source factor is one and the target factor is
$\eta(\operatorname{detId}(M,M_1))^{-n/2}$. For $h_1=h$ the target factor is one
and the source factor is $\eta(\operatorname{detId}(L,L_1))^{m/2}$.
See \nolinkurl{WeilInter.quot_bijective}, \nolinkurl{wd_congr},
\nolinkurl{om_iop} (\path{WeilRepresentationCommHeisenbergIntertwiner.lean:28--67,88--125})
and \nolinkurl{TrueTheta.exists_family_of_curve}
(\path{ExistsThetaVectorFamilyOfCurveDatum.lean:47--67,95--108}).

\paragraph{The scalar for rational movers.}\phantomsection\label{r14:eq-rational}
For $u=\det g$ rational unitary, $u\bar u=1$. Write $u=b/\bar b$,
using $b=1+u$ for $u\ne-1$ and $b=\sqrt{-d}$ otherwise. A Hecke
character of infinity weight $j_\tau$ and sign $\epsilon$ then gives
$\psi(b)^{\epsilon}=\prod_\tau\tau(u)^{-k_\tau}$ if
$2k_\tau=\epsilon j_\tau$. This follows by the principal-idele law and
$\tau(b)/\overline{\tau(b)}=\tau(u)$. Apply this to the exact Weil
vacuum multiplier with the above two weights. For a single rational
source mover it gives $\prod_\tau\tau(\det g)^{-m/2}$, and for a
single target mover it gives $\prod_\tau\tau(\det h)^{n/2}$.
See \nolinkurl{RatEqui.exists_root}, \nolinkurl{char_value},
\nolinkurl{det_root} (\path{OddEqOfEven.lean:26--71}) and
\nolinkurl{VacuumExact.vacuum_exact} (\path{WeilOperatorVacuumExactConstantEvenRank.lean:79--105}).

To compare the analytic pullback on $v$ with the same average, first
suppose $1-g$ is invertible. On each complex matrix space right
multiplication gives
$D_0=\det_{\mathbf C}(1-u_0^{-1})=\prod_\tau\tau(\det(1-g))^m$,
where $u_0^{-1}$ acts by right multiplication by $g_\tau$.
In the periodized Gaussian formula of Paragraph~\ref{par:r4-z43-iso}
(an earlier argument inside partial unit Z4.3, treated as a supplier by
review A)
the scalar is $|D_0|/D_0$ times the normalized Gauss phase. Here the
half-coordinate projection on the unchanged even-dimensional side splits
the finite quotient sum into two factors, interchanged by the coordinate
reversal, so the factors are equal. The auxiliary integral endomorphism
commuting with that projection conjugates the summand and permutes its
first factor. The first sum is real, and the full sum is its square.
The intertwiner and pullback have equal norms on $v\ne0$, so the
nonnegative phase is one. For $w=\tau(\det(1-g))$ the identity
$\det g\,\overline{\det(1-g)}=\det(1-g)$ follows by taking
determinants of $g(1-g^*)=-(1-g)$ and using even $n$. Hence
$|w^m|/w^m=\tau(\det g)^{-m/2}$. On the target side $1-h^*$ acts
on the left, repeated $n$ times. For
$w'=\tau(\det(1-h^*))$ one has
$\overline{w'}=\tau(\det h)w'$, giving
$|w'^n|/w'^n=\tau(\det h)^{n/2}$.
These two constants agree with the Weil multipliers, including their
signs. See \nolinkurl{SymConst.gauss_nonneg}, \nolinkurl{const_of_symm}
(\path{UnitaryThetaConstantEqInverseDeterminantPhase.lean:39--56,98--107,243--278}),
\nolinkurl{finsum_quotient_nonneg_of_swap_conj} (\path{FinsumQuotientNonnegOfSwapConj.lean:75--99}),
\nolinkurl{RatRow.pull_eq_smul_iop}
(\path{ThetaPullbackFirstGroupRationalPointEqSmul.lean:128--190}) and
\nolinkurl{RatEqui.pull_right_eq_det_smul}
(\path{ThetaPullbackSecondGroupRationalPointEqDetSmul.lean:150--180}).

Here is the integral endomorphism in the phase calculation. If $j$ clears
the denominators of $(1-u_0^{-1})^{-1}$ on the working period lattice,
choose $\ell,b\in\mathcal O_E$ with $\ell\bar\ell=-1+4jb$. To see the
congruence, factor $(4j)$ into prime powers. On a conjugate pair choose
residues $1,-1$. At a fixed prime $P$ the host's unramified conjugation
induces a nontrivial residue involution; no fixed prime lies above $2$.
The latter follows from $\bar w_0=1-w_0$ and $w_0(1-w_0)\in P$ above
$2$: primality and stability would put both $w_0$ and $1-w_0$ in $P$.
At such an odd fixed prime the norm from the finite residue field onto
the fixed field is surjective. Lift a solution modulo $P^a$, $a\ge1$,
by putting $d=z\bar z+1$ and $z_1=z(1+yd)$, where $2y=1$ modulo
$P$. Since $y+\bar y=1$ modulo $P$, the new error modulo $P^{a+1}$
is $d(1-y-\bar y)=0$. Chinese remaindering gives $\ell$ globally, and
conjugation of its norm equation shows $\bar b=b$.
Use entrywise multiplication by $\ell$ and $b$ on the unchanged side
of the matrix space. Both preserve the lattice and commute with the
half-coordinate projection and the moving unitary. Their pairing obeys
$H(\ell x,\ell y)=-H(x,y)+4jH(bx,y)$. Thus multiplication by $\ell$
is invertible modulo $j$, with inverse $-\bar\ell$. For the Cayley operator
$C=(1+u_0^{-1})(1-u_0^{-1})^{-1}/2$, write $jx=y-u_0^{-1}y$ with
$y,u_0^{-1}y$ in the lattice. Then $2jCx=y+u_0^{-1}y$ is integral.
Substitution in the phase
$\exp(\pi i(s\operatorname{Re}H(x,x)-\operatorname{Im}H(Cx,x)))$
reverses its sign in the exponent modulo $2\mathbf Z$. Indeed
$2s\operatorname{Re}H(bx,x)$ is integral by polarizing the integral
diagonal, and $2j\operatorname{Im}H(bCx,x)$ is integral by the
preceding identity. This proves the conjugate permutation of the first
half of the finite sum. See \nolinkurl{NegOneNorm.exists_norm_eq_neg_one},
\nolinkurl{exists_dvd_mul_involution_add_one}
(\path{ExistsDvdMulInvolutionAddOne.lean:12--79,81--127}),
\nolinkurl{NormResidue.exists_mul_involution_add_one_mem_pow}
(\path{ExistsMulInvolutionAddOneMemPow.lean:41--106}),
\nolinkurl{SymConst.gauss_nonneg}
(\path{UnitaryThetaConstantEqInverseDeterminantPhase.lean:140--192}),
\nolinkurl{RatRow.pull_eq_smul_iop} and \nolinkurl{RatEqui.pull_right_eq_det_smul}
(\path{ThetaPullbackFirstGroupRationalPointEqSmul.lean:197--219},
\path{ThetaPullbackSecondGroupRationalPointEqDetSmul.lean:196--234}).

The equality on $v$ extends to every rational Heisenberg translate of $v$
by covariance of both operators. At each finite level, character projection
on the finite period quotient expresses the section as a linear combination
of these translates: the standard fixed-character component is the
one-dimensional standard theta space, and translating by dual torsion
points gives the other components. Thus $W$ agrees with pullback on the
theta span. To include a rational matrix with eigenvalue $1$, factor
$g=(zg)(\bar z\,1)$ with $z\bar z=1$, $z\ne1$, and $z^{-1}$
outside the finite spectrum of $g$. Indeed for a nonzero imaginary
$\theta$, $(t+\theta)/(t-\theta)$ gives infinitely many distinct
norm-one scalars as $t$ varies in $\mathbf Q$; only finitely many are
excluded. See \nolinkurl{exists_normOne_scalar_avoiding_matrix_eigenvalues}
(\path{ExistsNormOneScalarAvoidingMatrixEigenvaluesStatement.lean:8--9}).
Both factors are admissible for the
preceding calculation. Multiply their operator identities; the same
factorization works for $h$. Finally the two groups act on different
sides of $z_\tau$. Multiplication of the two identities and
$W(\gamma_A g,\gamma'_A h)v=W(\gamma_A,\gamma'_A)W(g,h)v$
prove simultaneous covariance for $E$. See \nolinkurl{ThetaSpan.theta_span}
(\path{ThetaSpanRationalTranslatesVacuum.lean:143--176}),
\nolinkurl{RatExtend.om_eq_pull_of_vacuum}
(\path{OmEqPullOfVacuum.lean:207--227}),
\nolinkurl{RatEqui.exists_factor}, \nolinkurl{isPull_row}, \nolinkurl{isPull_col}
and the $a=a'=0$ specialization of \nolinkurl{oddEq_of_even}
(\path{OddEqOfEven.lean:82--109,117--180}; its full statement permits
arbitrary exponents satisfying the two weight equations),
and \nolinkurl{TrueTheta.exists_family_of_curve}
(\path{ExistsThetaVectorFamilyOfCurveDatum.lean:82--94}).

\paragraph{Change from the reference character to the prescribed one.}\phantomsection\label{r14:eq-normalize}
At $(S,S')$ the reference vector is nonzero. The tensor $B.\mathrm{thZero}$
belongs to this same theta line: its factors are entire and their
translation laws multiply. On a standard matrix period the quadratic
exponent is a sum of integral entry exponents, so its semicharacter is
exactly the product of theirs. Self-duality of the imaginary form gives
at most one independent theta vector, by the Fourier-coefficient
recursion in Paragraph~\ref{par:r4-z43-iso}. In applying that fact to
the actual period lattice, trace-dual families for the Hom lattice provide
an expansion of each vector; testing against the dual generators reduces
the alternating duality test to the elliptic period lattice. Hence
$B.\mathrm{thZero}=aE(S,S')$ with one $a$ depending only on the ranks
and fixed initial data. No division by $a$ is needed.
See \nolinkurl{WeilVac.thZero_mem}
(\path{ExistsWeilOThZeroTriv.lean:80--103}) and \nolinkurl{ThetaLine.theta_line},
\nolinkurl{of_dual_families}
(\path{ThetaSpaceSelfDualPairAtMostLine.lean:145--170,209--220,276--294}).

Write $C_\chi(L,M)=\chi(\operatorname{detId}(S,L))^{m/2}\chi(\operatorname{detId}(S',M))^{-n/2}$.
For a chain of self-dual lattices relative determinants multiply and are
norm-one ideals. Inverting the chain gives the reciprocal character value
because $\chi(1)=1$; in particular all these factors are nonzero. Transport
$E(S,S')$ first to $(L,S')$ and then to $(L,M)$ by the two single-lattice
laws. By complex linearity of each finite average,
\[
 e_\chi(L,M)=a\frac{C_\chi(L,M)}{C_\eta(L,M)}E(L,M).
\]
Composition of relative determinant ideals, including their norm-one
property for self-dual pairs, is the input \nolinkurl{RelDetIdeal.detId_normOne_mul}
(\path{DetIdNormOneMul.lean:156--166}). For $gL$ the relative determinant is multiplied by $(\det g)$; the forward
map is $g$, and the inverse inclusion is obtained with $g^*$.
Both $\chi$ and $\eta$ therefore acquire the same
$\prod_\tau\tau(\det g)^{-1}$. The target determinant has the
analogous multiplier. The ratio above is invariant under both moves,
and the covariance of $E$ proves the first displayed law. Changing
only $M$ to $M_1$ multiplies the ratio by
$[\chi(\operatorname{detId}(M,M_1))/\eta(\operatorname{detId}(M,M_1))]^{-n/2}$.
The denominator cancels against the reference target factor. Changing only
$L$ yields the same cancellation with exponent $m/2$. These are the
second and third laws on the same $B$ and $\chi$. See
\nolinkurl{chi_detId_mul}, \nolinkurl{chi_detId_ne_zero}, \nolinkurl{iop_mul}
(\path{MovedVacuumMemThetaFunctionsAndMoves.lean:65--87}),
\nolinkurl{ChiFacLaw.detId_image} (\path{DetIdLatGmap.lean:15--33}),
\nolinkurl{NatEqui.chi_smulLat}, \nolinkurl{sd_smulLat},
\nolinkurl{chiFac_ne_zero}, \nolinkurl{chiFac_right}, \nolinkurl{chiFac_left}
and \nolinkurl{NatEqui.eqP}
(\path{NormalisedFamilySatisfiesEquivarianceLaws.lean:17--57,59--117}).

\paragraph{The two uses on Z4.3.}\phantomsection\label{r14:eq-z43}
For the tied datum use $t.\mathrm{cc\_ok},t.\delta\mathrm{\_ok},
t.\chi\mathrm{\_ok}$ in the preceding proof. At an inert prime
$p\ge5$ with Jacobi symbol $(-d/p)=-1$, weak approximation in
\nolinkurl{KerInert.transOne} supplies $L=gM$, $L'=hM'$ where $M,M'$ are self-dual
of even ranks and locally standard at $p$.
Apply the first law to this very $M,M',g,h$. It changes the normalized
vacuum in the translated series to its pullback. The separate torsion
transport, coefficient algebraicity, standard-pair estimate and pullback
of the basis then apply. Lean takes the same conjunct from
\nolinkurl{eulerProduct_eq_of_curve_ideal}; the proof here supplies it
through the reference family. The reader for the class-sum argument is chosen separately. In
\nolinkurl{at_of_true_columns} the inert-series supplier
\nolinkurl{admissibleKernelSeries_inertPrime_nonConjugateFrobenius} takes
the adapted $t$ and $t.\mathrm{Hyp}$. The complex forms supplier
\nolinkurl{kernelClassSums_carry_chainComplexForms} takes
$t'$ with $\mathrm{ColSide}(t')$ and $\mathrm{FormsPin}(t')$.
These are exactly the existential witnesses of $\mathrm{Duty}(t).\mathrm{forms}$.
To compare characters, choose the witness actually constructed by
\nolinkurl{Pdv.duty_of_std}. Its standard-reader supplier
\nolinkurl{Pdv.stdKer_stdRd} rebuilds the record with the prescribed
reader and top chain and copies $\mathfrak c,\delta,B,\chi,c_D,O$ and
the entry function. Thus this explicit witness has the original
character. An arbitrary witness of the existential supplies its own
character for the same universal covariance law. See
\nolinkurl{Pdv.Duty} (\path{IsEntC.lean:39--43}),
\nolinkurl{at_of_true_columns} (\path{AtOfTrueColumns.lean:32--42}),
\nolinkurl{duty_of_std} (\path{ConditionsAtStandardKernelComparisonTrueColumns.lean:17--24})
and \nolinkurl{standardKernelData_with_standardBaseLevels}
(\path{StandardKernelDataWithStandardReadersStatement.lean:13--14},
\path{StandardKernelDataWithStandardBaseLevels.lean:15--16,31--42}).

For the complex stabilizer use a compatible
seam, the admitted seam lattices and a unitary $h$ stabilizing the target
lattice. Take $g=1$ and use the seam's unitary matrix of $h$. Both lattice
arguments stay fixed. Under reindexing the conjugate transpose on the
left is the inverse of the target action, so substitution gives
$e(L_hz)=e(z)$ in the reader's coordinates. The finite datum's own
multiplier is a separate input. See \nolinkurl{TieBase.base_data}
(\path{KernelEntryPeriodNeZeroAndLatticeLaws.lean:14--19}),
\nolinkurl{KerInert.transOne} (\path{InertCoprimeLemmas/IntAsm.lean:55--77}),
\nolinkurl{KerLaw.vac_leftE}, \nolinkurl{thC_left_stab}
(\path{ThCLeftStab.lean:368--397,432--436}).

\emph{Review boundary.} This derivation reconstructs both reference
transports, the exact scalar cancellation, and the two uses of the first
law. The construction of the finite action and the inert
vacuum calculation continue at pp.~\pageref{r14:weil-action}
and~\pageref{r14:weil-sign}, with their stated local inputs.
The reference coefficient field uses the algebraicity and conjugation of
averaged jets as in the vector-family argument on
p.~\pageref{r14:vec-fixed}. The review confirms the scalar cancellation and records the full
$\mathrm{EqP}$ and its first conjunct as conditional. For the first
conjunct the vacuum enters \nolinkurl{TrueTheta.exists_kappa} only
through membership in the standard theta space and nonvanishing
(\path{ExistsMovesOfSelfDual.lean:570--578},
\path{ExistsThetaVectorFamilyOfCurveDatum.lean:82--94}).
Taking an arbitrary such vector bypasses \nolinkurl{Vacuum.exists_vacuum},
its three Hasse-norm boundary declarations and the three CM coefficient
laws; the movers and the Weil identities remain inputs. The remaining Z4.3
application premises and its weighted kernel expansion are separate.

\paragraph{The finite Weil action on the translates of one vacuum.}\phantomsection\label{r14:weil-action}
Write $U(x)$ for Heisenberg translation, $A$ for the completed standard
lattice, and $v$ for its nonzero fixed vacuum. Its translates span the
all-level theta space by \nolinkurl{FinWeil.Vacuum.vacuum_fixed_span}
(\path{AdelicWeilVacuumFixedAndTranslatesSpan.lean:57--90}). The passage
to the adelic model uses \nolinkurl{FinWeil.Model.exists_isModel}
(\path{ExistsAdelicWeilRepresentationModel.lean:67--129}), and its
identification with the analytic average uses \nolinkurl{FinWeil.iop_eq_sumOp}
(\path{ThetaIntertwinerEqNormalisedAdelicSum.lean:60--127}).
For a unitary pair $g$, put $A_g=gA$ and prescribe
\[
  L_g U(x)v=U(gx)v_g,\qquad
  v_g=\lambda([A_g:A_g\cap A])^{-1}
               \sum_{b\in A_g/(A_g\cap A)}U(b)v .
\]
The sum uses representatives. Changing a representative by the intersection
does not alter the summand by the fixed-vector and isotropy laws, and
translation by $A_g$ fixes the sum. Here is why the prescription respects
relations. For $a\in A$, the vectors $U(x)v$ have eigencharacter
$a\mapsto\gamma(a,x)$ under $U(a)$. Distinct cosets modulo $A$ give
distinct characters by self-duality. A finite relation therefore splits
by these characters: acting by $U(a)-\gamma(a,x_0)$ removes a chosen
term, and induction separates any two different characters. Within one
coset the translates are proportional by the Heisenberg multiplication
law. The primed vectors $U(gx)v_g$ obey exactly the same proportionalities
and character distinctions under $U(ga)$, because $g$ preserves
$\gamma$. Thus $L_g$ exists on the span. On a translate, composition
reduces to sending the representatives for $A_k/(A_k\cap A)$ to their
images under $g$, then composing the two normalized averages. The
composition formula for self-dual lattices gives $L_g L_k=L_{gk}$.
Its induction chooses an intermediate adjacent lattice, reducing the
relative index by a residue-field cardinal greater than one. The
adjacent-lattice composition formula is a separate input of that
induction. The normalization uses $\lambda(p)=\sqrt p$ off the
inert primes and $\sqrt p\exp(\pi i/f_p)$ at an inert prime of residue
degree $f_p$, so $\lambda(p^{f_p})=-\sqrt{p^{f_p}}$ there. On split-supported
moves the relative index has no inert prime divisor; the normalized
adelic sum becomes the positive-square-root theta average. Multiplying
$L_g$ by the two determinant characters gives the split vacuum scalar
above. See \nolinkurl{FinWeil.exists_weil_rep}
(\path{ExistsWeilRepresentationHeisenberg.lean:81--120,135--192,198--278}),
\nolinkurl{FinWeil.normOp_comp}
(\path{NormalizedHeisenbergIntertwinerComp.lean:14--44}),
\nolinkurl{FinWeil.exists_lam}, \nolinkurl{FinWeil.exists_goodNorm}
(\path{ExistsGoodIndexNormalizationNegativeInert.lean:34--73,92--125}) and
\nolinkurl{FinWeil.exists_weil_of_vacuum} (\path{ExistsWeilOfVacuum.lean:22--42,71--83}).

\paragraph{The sign of an inert neighbour.}\phantomsection\label{r14:weil-sign}
Let $k$ be the residue field at the conjugation-stable maximal ideal,
with nonidentity involution $s$ and fixed field $F$. In the local
triangle of rank $r$, the two quotient coordinates are in $k^r$.
The condition that the sum of their representatives lies in the first
lattice identifies the coordinates on the diagonal. The off-diagonal
pairings between different indices are trivial, leaving one factor
$\sum_{t\in k}\psi_i(t s(t))$ for each coordinate. Here $\psi_i$
is an additive character trivial on the defining ideal. Its restriction
to $F$ is nontrivial: triviality on all $a+s(a)$ would put the first
generator for that coordinate in the third lattice by self-duality,
and the two quotient-coordinate kernel tests would then imply
$1$ lies in the defining ideal. To evaluate the factor, write
$q=|F|$. The nonidentity involution on a finite field has order two,
so $|k|=q^2$, $s(t)=t^q$, and the norm $t\mapsto t^{q+1}$ has
$q+1$ preimages on each nonzero element of $F$ by cyclicity of $k^\times$.
Its zero fibre has one element. Character orthogonality gives
\[
 \sum_{t\in k}\psi_i(t s(t))
     =1+(q+1)\sum_{u\in F^\times}\psi_i(u)=-q .
\]
The triangle multiplier is the product of these factors times a
positive real normalization. Its eighth power is one, hence the
multiplier is $(-1)^r$. For a single neighbour $B\to C$, choose the
local transvection taking $B$ to $C$ and fixing a third neighbouring
lattice $A$. Its transported vector at $A$ is fixed because two
inert-supported movers with the same lattice differ by an integral
stabilizer. Express the vector at $B$ as a multiple of the average
from $A$ using the theta-line property. Transport gives the same multiple
for $C$, and the triangle just evaluated supplies the sign for the
step $B\to C$. A neighbour towards the standard lattice lowers
$\ell(S,L)=\operatorname{length}(S/(S\cap L))$ by one. The intermediate
distributive triple has multiplier one, so induction gives
$(-1)^{m\ell(S,L)+n\ell(S',M)}$ for both sides together, with ranks
$n,m$ respectively. It is one for the two even ranks here. This
derivation uses the local transvection, diagonal neighbour coordinates,
and the positive-normalization/eighth-root property of the triple as
inputs. See \nolinkurl{NormGauss.exists_norm_eq}, \nolinkurl{NormGauss.sum_norm_neg}
(\path{FinsumNormNeg.lean:8--55,57--183}),
\nolinkurl{TriangleSign.mu_eq_sign}
(\path{ThetaTripleConstantEqNegOnePow.lean:98--161,173--221}), and
\nolinkurl{InertSign.exists_move}, \nolinkurl{InertSign.core}, \nolinkurl{InertSign.vacuum_inert}
(\path{WeilRepresentationVacuumSignInertPlaces.lean:89--119,144--196,197--236,275--327}).

% blueprint/ch_P_r14.tex lines 1060 through 1110; r14:base-assembly
\subsection{Assembly and downstream use}\label{r14:base-assembly}
For the given $t$, use the nonzero period of Section~\ref{r14:base} in
\nolinkurl{Base.latAll}. This proves both the single-pair lattice laws and
the two ordered indices between any admitted pairs. Specialize
\nolinkurl{ColNorms.familyVectors} to $(D,\mathfrak c,\delta,B,\chi)$ with
the three host and finite-support arguments listed in
Section~\ref{r14:vec}. This gives membership, nonvanishing and one number
field for all coefficient indices at each pair. Section~\ref{r14:equivariance}
specializes \nolinkurl{Mem.NatEqui.eqP} to the same three correctness
fields of $t$. Its \emph{single scalar for fixed ranks} is determined at
the standard pair and transported to every pair. The scalar used in the
coefficient-field descent is allowed to depend on the pair. Thus both
universal statements concern $B.e^{\mathrm{nat}}(\chi)$ with their declared
quantifiers. Taking the indicated conjunction proves the proposed
implication of \nolinkurl{GlobalTheory.TieBase.base_data}
(\path{SatakeParameters/KernelEntryPeriodNeZeroAndLatticeLaws.lean:14--19}).
Review A confirmed as hand-proved the conjunction $\delta\ne0\ \land B.\mathrm{LatAll}$ in this partial unit. The assembly and character
factors match the root. The proofs of $\mathrm{VecP}$ and
$\mathrm{EqP}$ still depend on the suppliers specified above. Neither $t.\mathrm{Hyp}$ nor a second character
is a premise of this assembly.

\paragraph{Where these conclusions enter the next units.}
Z4.2.09 uses the four conclusions to specialize the split-site law with
the pair, position, frame movement and away-from-place laws on the same
site (\nolinkurl{TwV.site_law_xiIn},
\path{SatakeParameters/KernelSiteLawAndFixedCoordinateStabiliserLaw.lean:20--45}).
The pullback for Z4.2.10 and .11 receives them through
\nolinkurl{kernelEntry_period_ne_zero_and_lattice_laws} in
\nolinkurl{SPull.kerTh_eq_pull}
(\path{SatakeParameters/ThetaKernelFunctionAdelicPointsEqPullbackClassKe.lean:65--79}).
At Z4.2.12 the site supplies the dictionary, pair, position, move,
away and scalar laws. One embedding extending that site's $\lambda$ serves
both the zero-step and positive-step frame comparisons.
See \nolinkurl{admissible_at_each_site}, \nolinkurl{site_cmpLite_siteAdm}
(\path{SatakeParameters/JointSiteChainComparedCarriesAdmittedSystem.lean:29--49,62--83}).
The split-place interpolation, independent class-kernel comparisons and
assembly retain their own additional conclusions and hypotheses.
The twisted-prime law Z4.2.02 is a further open unit.

\paragraph{Which dependency size is being counted.}
At the pinned source, traversing type and value edges of
\nolinkurl{TieBase.base_data} until the first literature declaration gives
9,400 core declarations and eight literature leaves. The core contributes
60,733 inclusive declaration lines. Counting only the 3,766 core
theorems with valid ranges gives 51,444 lines; removing the 176 singleton
ranges gives 51,268. These are three measures of the same root. Its
entire 60,733-line cone belongs to each of Z4.2.09--12. The reference-family
route \nolinkurl{NatEqui.eqP} accounts for 50,356 core lines inside it.
This last number measures dependencies, not how many lines have received
independent mathematical review (\path{blueprint/status.json},
\path{out/R14_pending.json}).

% Parent source SHA-256: 726efab7fd295a2dce4fc43b2a38c6ed087511bfb339c03e564ed0607ef7693c
% blueprint/ch_P_r15.tex lines 1092 through 1183; r15:weil
\paragraph{Constructing an operator with prescribed characters.}
\phantomsection\label{r15:weil}
For even $n,m$, every $v\ne0$ and every pair of Hecke characters
$\psi,\psi'$ of ranks $n,m$ admit a $W$ with
$(W.\mathrm{chV},W.\mathrm{chV}',W.\mathrm{ep},W.\mathrm{ep}')
=(\psi,\psi',1,1)$.
The construction works for arbitrary ranks in
\nolinkurl{FinWeil.exists_weil_of_vacuum}, which specializes to
\nolinkurl{FinWeil.exists_weil}
(\path{FoundationProofs/ExistsWeilOfVacuum.lean:22--83}).

Use the adelic Heisenberg model on the finite-level space
$\mathcal H=\mathrm{TT}$, with operators $U(x)$, commutator pairing
$\gamma$, self-dual standard lattice $A$ and action
$\rho(g,h)x=hxg^{-1}$. The model supplies agreement on rational points
and, on each vector, constancy on a sufficiently small open coset.
The standard theta law and density make $v$ fixed by $A$. Restriction
to rational periods and the one-dimensional theta bound show that the
$A$-fixed vectors are $\mathbf C v$. A vector fixed by $lA$ splits
under the finite abelian action of $A/lA$ into character components.
Self-duality pairs this group perfectly with $l^{-1}A/A$; applying an
inverse Heisenberg translate turns each component into an $A$-fixed
vector. Hence the $U(x)v$ span $\mathcal H$.
These uses of \nolinkurl{Model.exists_isModel},
\nolinkurl{Vacuum.vacuum_fixed_span}, and \nolinkurl{span_translates_vacuum}
are in \path{FoundationProofs/ExistsAdelicWeilRepresentationModel.lean:67--129},
\path{FoundationProofs/AdelicWeilVacuumFixedAndTranslatesSpan.lean:36--90},
and \path{FoundationProofs/SpanHeisenbergTranslatesVacuumEqTop.lean:10--76}.

For self-dual commensurable $A,B$ and $u$ fixed by $A$ use
$T_{AB}u=\lambda([B:A\cap B])^{-1}\sum_{b\in B/(A\cap B)}U(b)u$.
The good normalization $\lambda$ equals the positive square root on
indices without inert factors. On a simple quotient by a
conjugation-stable prime $P$, its value is $-\sqrt{|k_P|}$.
The host condition $D.\mathrm{host.der.unramEF}$ excludes ramification
in the relative extension, so the nonsplit primes here are inert
(\nolinkurl{FinWeil.exists_goodNorm},
\path{FoundationProofs/ExistsGoodIndexNormalizationNegativeInert.lean:92--125}).
Each sum is $B$-fixed by reindexing. The composition rule
$T_{BC}T_{AB}u=T_{AC}u$ follows from the adjacent-lattice sum identity
by induction on $[B:A\cap B]$: take an adjacent $A'$ towards $B$
with $d_{AB}=d_{AA'}d_{A'B}$ and $d_{AA'}>1$, where $d_{AB}=[B:A\cap B]$.
Apply the adjacent identity first with target $B$, use induction from
$A'$ with target $C$, then apply the adjacent identity with target $C$.
At index one, self-duality implies $B=A$ and $T_{AA}u=u$. The local
input for the equal-intersection adjacent case is the signed Gauss
factor; the construction also uses existence of $A'$.
See \nolinkurl{exists_goodNorm}, \nolinkurl{sumOp_comp_of_adj_of_inf_eq},
\nolinkurl{normOp_comp} and \nolinkurl{exists_adj_selfDual_step}
(\path{FoundationProofs/ExistsGoodIndexNormalizationNegativeInert.lean:92--125},
\path{FoundationProofs/HeisenbergSumCompOfAdjacentOfInfEq.lean:22--192},
\path{FoundationProofs/NormalizedHeisenbergIntertwinerComp.lean:14--45},
\path{FoundationProofs/ExistsAdjacentSelfDualLatticeToward.lean:65--104}).

For $s=(g,h)$ put $A_s=\rho(s)A$, $v_s=T_{A A_s}v$, and define
$L_sU(x)v=U(\rho(s)x)v_s$. To see that the formula respects relations,
separate the $A$-eigencharacters of the translates in a finite
relation. Distinct cosets modulo $A$ give distinct characters by
self-duality. Inside one coset the Weyl product expresses every
translate as a scalar multiple of one nonzero translate. The same
scalar occurs after transport because $v_s$ is $A_s$-fixed and $\rho$
preserves the pairing and the multiplication phase. Thus $L_s$ extends
linearly and intertwines $U$ with $U\circ\rho(s)$. Transporting the
summation representatives by $\rho(s)$ gives
$L_sL_t v=T_{A_s A_{st}}T_{A A_s}v=v_{st}$; intertwining extends
this equality and $L_1=1$ to the spanning translates
(\nolinkurl{exists_weil_rep},
\path{FoundationProofs/ExistsWeilRepresentationHeisenberg.lean:105--197,198--278}).

Twist $L_s$ by $c_\psi(h)c_{\psi'}(g)$, with
$c_\eta(g)=\eta(1-\varepsilon+\varepsilon\det g)$, where
$\varepsilon$ selects one component above each split prime. This is a
character, trivial on inert-supported movers. If
$x=\det(g)\bar x$ and $x=1$ at inert primes, the quotient of $x$ by
$1-\varepsilon+\varepsilon\det g$ is $b\bar b$, where
$b=1-\bar\varepsilon+\bar\varepsilon x$. The norm law gives
$c_\eta(g)=\eta(x)$. Extend the twisted operator by zero outside
$\mathcal H$. On $\mathcal H$ it preserves the space and satisfies
additivity, homogeneity, identity and composition. Constancy at
$\rho(s)x$ supplies the modulus for the rational lift law. For
integral inert-supported movers $A_s\subset A$, so the sum on $v$
has one term and the twist is one. At split-supported movers the index
has no inert factor and the sum equals the positive-normalized analytic
intertwiner: the comparison uses rational representatives of
$A_s/(A_s\cap A)$ and agreement of $U$ with rational Weyl operators.
The multiplier on $v$ is $\psi(x')\psi'(x)$. This verifies the three
remaining operator laws. The index and coset comparisons used here
are \nolinkurl{idx_map_actA} and \nolinkurl{iop_eq_sumOp}; the twist
is \nolinkurl{SplitCharacter.exists_split_character}
(\path{FoundationProofs/HeisenbergRelativeIndexTranslateCompletedLattice.lean:13--32,48--108},
\path{FoundationProofs/ThetaIntertwinerEqNormalisedAdelicSum.lean:60--127},
\path{FoundationProofs/ExistsSplitCharacterUnitaryGroup.lean:71--116}).

% blueprint/ch_P_r15.tex lines 1185 through 1260; r15:vacuum-inert
\paragraph{The sign at inert places.}\phantomsection\label{r15:vacuum-inert}
Let $nm$ be even and let $W$ be any Weil operator on $v$; its field
$W.\mathrm{nz}$ gives $v\ne0$. For inert-supported movers $g,h$
taking $S,S'$ to $L,L'$, set $A_f=\mathrm{La}(L,L')$ and
$\ell(B,C)=\mathrm{length}_{\mathcal O_E}(B/(B\cap C))$. Then
\[
 W(g,h)v=(-1)^{n\ell(S',L')+m\ell(S,L)} I_{A_0 A_f}v.
\]
In \nolinkurl{InertSign.vacuum_inert} the analytic input is that each
moved theta space is at most a line, with mutually inverse normalized
intertwiners, equivariant for $W$ on moved pairs
(\path{FoundationProofs/WeilRepresentationVacuumSignInertPlaces.lean:254--327};
\nolinkurl{ThetaLine.theta_line}, \nolinkurl{WeilInter.om_iop},
\path{FoundationProofs/ThetaSpaceSelfDualPairAtMostLine.lean:276--294},
\path{FoundationProofs/WeilRepresentationCommHeisenbergIntertwiner.lean:88--125}).

The moved vector belongs to that line by
\nolinkurl{WeilTheta.om_mem_Th_moved}
(\path{WeilRepresentationThetaMovedLattice.lean:21--43}).
\nolinkurl{T.kernelLattices_form_selfDual_triples} provides the
triple multiplier and \nolinkurl{T.exists_traceOne_integer_conj_diff_prime_to}
supplies a trace-one integer whose conjugate difference is nonzero
modulo the chosen prime (\path{KernelLatticesFormSelfDualTriplesStatement.lean:7--8},
\path{ExistsTraceOneIntegerConjDiffPrimeToStatement.lean:7--8}).

Move one coordinate, keeping opposite rank $r$ fixed. Write $E_B$ for
the image of $v$ under an inert mover to $B$. Two such movers differ
on the right by an integral inert stabilizer; hence $E_B$ is
well-defined, and the inverse action makes it nonzero. For a nonzero
length choose a simple submodule of $S/(S\cap B)$ with generator
$x$ lifted to $S$; write $P$ for its annihilator. Then $x\notin B$,
$Px\subset B$, and $\langle x,x\rangle$ is integral. The prime $q$
under $P$ is inert: otherwise $S$ and $B$ agree at $q$, so some
$s$ prime to $q$ carries $x$ into $B$; B\'ezout applied to $s$ and
$q$ would give $x\in B$. The nearer neighbour
$C=\{b\in B:\langle b,x\rangle\in\mathcal O_E\}+\mathcal O_E x$
reduces the length by one. The local transvection input supplies a
third reachable lattice $A$ and an inert mover from $B$ to $C$ fixing
$A$. For the theta lattices associated to these three arithmetic
lattices write $I_{AB}$, etc., for their intertwiners
(\nolinkurl{NearerNeighbour.exists_nearer}, \nolinkurl{InertSign.exists_move},
\nolinkurl{InertSign.core},
\path{FoundationProofs/CardParitySignLengthOfEq.lean:68--106},
\path{FoundationProofs/WeilRepresentationVacuumSignInertPlaces.lean:89--119,126--236}).

For this triangle the row or column coordinate theorem over $k=\mathcal O_E/P$
identifies the contributing pairs of cosets with equal residue coordinates.
The fixed field $k_0$ has order $Q$, with $[k:k_0]=2$. The
trivial-character component of the Heisenberg sums is a product of $r$
sums of $a\mapsto\chi_i(a\bar a)$ on $k$. Here each $\chi_i$ restricts
nontrivially to $k_0$: the diagonal-pairing lemma proves that triviality
on traces would make a generator commute with the whole opposite
lattice, contradicting self-duality and its nonzero residue. The norm
has $Q+1$ preimages over every nonzero element, so each sum is
$1+(Q+1)(-1)=-Q$. As the positive-normalized intertwiners have a
unit-modulus triple scalar, that scalar is $(-1)^r$. For a between-lattice
or mixed-coordinate distributive triple each contributing phase is one
by splitting the middle vector into its two intersections; the scalar is
then $+1$. The lower inputs here are the diagonal coordinates, their
nondegeneracy, and distributivity, as in
\nolinkurl{TriangleSign.mu_eq_sign}, \nolinkurl{DistLa.dist_La},
\nolinkurl{DistTriple.const_eq_one}
(\path{FoundationProofs/ThetaTripleConstantEqNegOnePow.lean:98--221},
\path{FoundationProofs/ThetaLatticeTriplesDistributive.lean:103--156},
\path{FoundationProofs/ThetaTripleConstantEqOneDistributive.lean:10--45}).

Write $E_B=c I_{AB}E_A$ using the line and inverse intertwiner.
The mover above fixes $E_A$, so equivariance gives $E_C=c I_{AC}E_A$.
The triangle then gives $I_{BC}E_B=(-1)^r E_C$. The reverse
intertwiner, the length decrement, and the between-lattice identity
give $E_B=(-1)^{r\ell(S,B)}I_{SB}v$ by induction, starting from the
identity mover. Apply this successively in the two coordinates.
Equivariance, the representation law and the mixed-coordinate scalar
one give the claimed exponent. Empty coordinate sets give sign $+1$
(\nolinkurl{InertSign.core}, \nolinkurl{InertSign.vacuum_inert},
\path{FoundationProofs/WeilRepresentationVacuumSignInertPlaces.lean:126--236,293--327}).

% blueprint/ch_P_r15.tex lines 1262 through 1328; r15:vacuum-exact
\paragraph{The exact multiplier.}\phantomsection\label{r15:vacuum-exact}
Now let $n,m$ each be even. For an arbitrary $W$ on $v$, write
$\psi=W.\mathrm{chV}$, $\psi'=W.\mathrm{chV}'$ and
$\epsilon=W.\mathrm{ep}$, $\epsilon'=W.\mathrm{ep}'$. For any
movers $g,h$ to $L,L'$ and ideles satisfying $x=\det(g)\bar x$,
$x'=\det(h)\bar x'$, one has
\[
 W(g,h)v=\psi(x')^{\epsilon}\psi'(x)^{\epsilon'} I_{A_0 A_f}v .
\]
First, $G=I_{A_0 A_f}v\ne0$ by applying the reverse intertwiner.
By \nolinkurl{T.weilOperator_conj_heisenberg_moved_lattice}
(\path{WeilOperatorConjHeisenbergMovedLatticeStatement.lean:8--9}),
every period in $A_f$ fixes
$W(g,h)v$: conjugating by the inverse mover gives a period of $A_0$
on $v$. This calculation takes place at finite level, since the level
of $W(g,h)v$ times $[A_0:A_0\cap A_f]$ also works after a shift by
an $A_f$ period. Thus $W(g,h)v=cG$ on the moved theta line
(\nolinkurl{VacuumExact.exists_scalar}, \nolinkurl{VacuumExact.vacuum_exact},
\path{FoundationProofs/WeilOperatorVacuumExactConstantEvenRank.lean:17--52,79--105}).

Let $e=\varepsilon+\bar\varepsilon$ project onto the non-inert
components. Cutting a mover into $eg+(1-e)1$ and $(1-e)g+e1$ gives
commuting split and inert parts with product $g$; likewise for $h$.
Let $A_s,A_i$ be their respective period lattices. The indices
$d_j=[A_j:A_j\cap A_0]$ satisfy $d_f=d_s d_i$: at inert primes the
split lattice agrees with $A_0$ and the full lattice with $A_i$,
and conversely at the other primes. This applies the local period
comparison in \nolinkurl{NormSum.La_agree_of_trivAt} and the
intersection-index identity \nolinkurl{VacuumScalar.index_mul}
(\path{FoundationProofs/ThetaLatticeAgreeOfTrivialAtPrime.lean:74--112},
\path{FoundationProofs/OddnessConditionHolds.lean:63--91,434--456}).

Let $c_s,c_i$ be the scalars for the parts. Write $H$ for the raw
full sum, so $Hv=\sqrt{d_f}G$. Transport the inert representatives
through the split mover into $A_f$ and adjust the split representatives
by periods in $A_s\cap A_0$ to put them in $A_s\cap A_f$. The first
step uses Weyl equivariance and the second uses local agreement and the
fact that the adjustment fixes $v$. Their raw sums $H_i,H_s$ commute
with $H$: shifts in $A_f$ have exact addition on $A_f$, as their
multiplication phase is one when tested on its nonzero fixed vector
$G$. Thus $HG=d_fG$, $H_iG=d_iG$, $H_sG=d_sG$. The representation
law gives
$cG=(c_sc_i)/\sqrt{d_id_s}\,H_iH_s v$. Apply $H$ and cancel $G$:
$c d_f=(c_sc_i)\sqrt{d_f}\,d_id_s/\sqrt{d_id_s}$, so $c=c_s c_i$
(\nolinkurl{VacuumScalar.scalar_eq}, \nolinkurl{VacuumScalar.scalar_mul},
\path{FoundationProofs/OddnessConditionHolds.lean:209--251,457--493}).
For the inert part both ranks are even, so its length sign is one.
The split determinant lifts $y=ex+1-e$, $y'=ex'+1-e$ are one at
inert primes. The ratios $x/y,x'/y'$ are supported there, and
\nolinkurl{InertTrivial.character_trivial_at_inert} therefore preserves
both even-rank character values. The field $W.\mathrm{split}$ gives
$c_s=\psi(x')^{\epsilon}\psi'(x)^{\epsilon'}$, proving the formula
(\nolinkurl{VacuumExact.split_part},
\path{FoundationProofs/WeilOperatorVacuumExactConstantEvenRank.lean:54--105},
\path{FoundationProofs/CharacterTrivialAtInertPrimes.lean:255--309}).
The three proofs use a given nonzero vector in the standard theta
space. In the two base-free applications (rows 18 and 22 of the
R14 review), $\delta\ne0$ is supplied by the curve
discriminant. Write the unscaled ideal lattice as $a(\mathbf Z\tau+
\mathbf Z)$ with $a\ne0$ and $\Im\tau>0$. Its discriminant is
$(2\pi/a)^{12}\Delta(\tau)$; the twelfth-root law equates this
nonzero value to $\delta^{12}$. This proves the needed nonvanishing
using just the curve datum
(\nolinkurl{CurveLattice.discr_of_tau}, \nolinkurl{CurveLattice.lattice_of_curve},
\path{FoundationProofs/IdealPeriodLatticeOfCurveDatum.lean:52--62,136--157}).
The coefficient field, CM conjugation and the local forms of
the full \nolinkurl{Vacuum.exists_vacuum} remain separate supplies.

% blueprint/ch_P_r15.tex lines 1330 through 1371; r15:jets
\paragraph{Translated coefficients for the prescribed natural vector.}
\phantomsection\label{r15:jets}
For a curve datum, base $B$, admissible character and admitted pair,
if a positive multiple of $x$ belongs to the period lattice then every
normalized Taylor coefficient of $U_x B.e^{\mathrm{nat}}$ is algebraic.
Section~\ref{r15:translated} checks the four-copy construction of
the standard tensor and the finite intertwiners, giving algebraic
coefficients throughout its rational period saturation. Explicitly,
$B.e^{\mathrm{nat}}=\mathrm{chiFac}\,B.\mathrm{thMoved}$; the coefficient
$\mathrm{chiFac}$ is algebraic at an admitted pair, and clearing the
positive multiple puts $x$ in that saturation. Linearity proves the
assertion in \nolinkurl{ThLoc.algTAll}
(\path{SatakeParameters/AlgTAll.lean:6},
\path{GaloisGeometry/AlgTAll.lean:12--26}). For the broader input
\nolinkurl{AlgTranslate.coefIn_translate}, start with any algebraic-jet
theta vector on two self-dual lattices. Frame restriction writes it as
a finite sum $\sum_i c_i f_i$ of zero-padded restrictions of torsion
translates of a product theta. The CM translation law on that product
gives $a_\alpha(f_i)\in k=\overline{\mathbf Q}$. Choose a $k$-linear
retraction $\pi:\mathbf C\to k$ of the inclusion. For every $\alpha$,
\[
 a_\alpha(e)=\pi\!\left(\sum_i c_i a_\alpha(f_i)\right)
            =\sum_i\pi(c_i)a_\alpha(f_i).
\]
Uniqueness from the Taylor series gives the same function with
coefficients $\pi(c_i)$. Padding the shift puts $Rjx$ (in padded
coordinates) in the product lattice, where $jx$ was an original
period and $R>0$ is the lattice comparison index multiple. Translation
and normalized Taylor coefficients commute with zero-padding.
The broader argument needs frame restriction, CM product translation
and the lattice comparison. In Lean, \nolinkurl{JetAlgebraic.ta_thMoved} uses
\nolinkurl{AlgTranslate.coefIn_translate}, which accounts for 10,897
theorem lines outside $D7$. Section~\ref{r15:translated} substitutes
the four-variable identity and the genus-one torsion translation for
that formal route, bypassing both this declaration and
\nolinkurl{cmTheta_taylorCoeff_algebraic_and_conjugate}. Its two
upstream arguments have independent locations in
Sections~\ref{r16:ktheta} and~\ref{r16:stdok}
(\path{VacuumEntireAndNormalizedJetsAlgebraic.lean:311--316}).
For the finite-span calculation see
(\nolinkurl{AlgTranslate.exists_alg_comb}, \nolinkurl{AlgTranslate.coefIn_translate},
\path{SatakeParameters/CoefInTranslate.lean:28--51,53--85}).

% Parent source SHA-256: 9f5cb0d34db7a3512d516aefedd7111b0e4c8718613c492875db00e44bd18a5f
% blueprint/ch_P_r8_base.tex lines 1 through 49; r8:Z4.2.13
\subsection{Z4.2.13. The period and theta family on the tied base}\label{r8:Z4.2.13}
Let \(t\) be a tied kernel entry for the fixed host, curve ideal, discriminant root, extra comparison data and lattice system. Write \(\delta=t.\delta\), \(B=t.B\), \(\chi=t.\chi\). The required conclusion is
\[
 \delta\ne0,\qquad B.\mathrm{LatAll},\qquad B.\mathrm{VecP}(\chi),\qquad B.\mathrm{EqP}(\chi).
\]
The first lattice clause gives the period-lattice conditions at every admitted pair. The second gives finite relative indices in both directions between the period lattices of any two admitted pairs of the same ranks. The vector clause requires a nonzero theta vector and a single number field for all of its Taylor coefficients at each pair. Equivariance comprises the simultaneous rational unitary change of basis and the two separate lattice changes with their determinant-character factors. All clauses concern the same normalized family \(B.eNat(\chi)\). See \nolinkurl{GlobalTheory.TieBase.base_data} (\path{SatakeParameters/KernelEntryPeriodNeZeroAndLatticeLaws.lean:14--19}), \nolinkurl{Base.LatAll}, \nolinkurl{Base.VecP}, \nolinkurl{Base.EqP} (\path{KernelLattices/LateralTheory.lean:173--175,185--189,207--219}) and \nolinkurl{theta_vector_family_iff_parts} (\path{GlobalTheory/GL3/PFamily.lean:13--19,34--36}).

\paragraph{Nonvanishing of the period.}
The two periods of the lattice in \(B\) are linearly independent. If \(\delta=0\), the lattice identified by the base with \(\mathrm{lamOf}(D,\mathfrak c,\delta)\) is zero. Its first period belongs to it, so is zero, contradicting independence. This proves the first clause for every tied base. See \nolinkurl{base_linkage_nonempty} (\path{GlobalTheory/GL3/PFamily.lean:80--86}).

\paragraph{One coefficient field for the entire germ.}
Fix an admitted pair. Suppose the reference family supplies a nonzero entire vector \(e\) with every Taylor coefficient in a single number field \(M\), and that the common theta-line comparison gives \(B.eNat(\chi)=a e\). Suppose also that all Taylor coefficients of the left side are algebraic, as obtained by the source from the translated-jet theorem at shift zero. Some Taylor coefficient \(e_\alpha\) is nonzero, since otherwise the Taylor expansion of the entire function would vanish identically. The linearity of the normalized Taylor coefficient gives
\[
 a=\frac{[z^\alpha] B.eNat(\chi)}{[z^\alpha]e}\in\overline{\mathbf Q}.
\]
Thus \(M(a)\) is finite over \(\mathbf Q\). Every coefficient of the normalized vector lies in this same field. The construction allows the field to depend on the pair. This proves the coefficient-field descent once the three stated inputs are supplied (\nolinkurl{GlobalTheory.ColNorms.familyVectors}, \path{DivisorProofs/NormalisedFamilyNonzeroThetaVectorsNumberField.lean:15--50}).

The application of \nolinkurl{ColNorms.familyVectors} in
\nolinkurl{TieBase.base_data} passes three further theorems
(\path{SatakeParameters/KernelEntryPeriodNeZeroAndLatticeLaws.lean:16--19};
\nolinkurl{FamilyVectors}, \path{DivisorProofs/CosetRefinement.lean:84--86}).
\nolinkurl{GL3Theta.cm_field_facts_hold} gives the host field conditions
\(\mathrm{adm1}(D)\), including the finite-place ramification and splitting
conditions. \nolinkurl{GL3Theta.induced_cm_type} gives
\(\mathrm{embT}_\tau(a)=\varphi_0(a)\) for every \(\tau\) and \(a\in K\)
viewed in \(E\) (\path{UnitTraceSumBounds/InducedCmType.lean:10--17};
\path{GlobalSchurBounds/GL3Theta.lean:828--831};
\path{GlobalSchurBounds/TaylorLogMetric.lean:206--215}).
\nolinkurl{ExponentsFinite.yFinite}, applied to the base's
choice \(k=B.\mathrm{kc}\), proves that the rational primes split in \(K\)
with nonzero relative-determinant exponent \(k.y_P(p)\) form a finite set.
It clears denominators for the lattice map and its inverse to bound those
primes (\nolinkurl{KapChoice.YFinite}, \path{KernelLattices/LateralTheory.lean:107--113};
\path{DivisorProofs/KernelConstantExponentsFinite.lean:85--126}).
These are explicit arguments to the family theorem. The quoted
coefficient-field deduction begins after choosing the reference family.

\paragraph{Expert lemmas still needed.}
\emph{[Proved in Lean, not reproduced: full lattice and family construction.]}
For the given curve ideal, its discriminant root and this base, establish self-duality for the imaginary part of the period pairing, its integral semicharacter law, and finite index in both directions. The scalar nonvanishing argument alone supplies none of these comparisons (\nolinkurl{Base.latAll}, \path{KernelLocalFactors/IsAdmissibleFlagPairOfUnitriangularLeftAutomorph.lean:623--644}). For these same admissible even-rank pairs, construct the nonzero reference family, prove that its theta space is a line, and identify the natural normalized vector with a nonzero member. The descent above uses these data before it tests a Taylor coefficient (\nolinkurl{familyVectors}, \path{DivisorProofs/NormalisedFamilyNonzeroThetaVectorsNumberField.lean:15--32}).

For a reference character \(\chi'\), write \(C_\chi(L,L')=\mathrm{Base.chiFac}(\chi,N,N',L,L')\).
With ranks fixed the required comparison uses one scalar \(a\) for all admitted \((L,L')\):
\[
 e^{\rm nat}_\chi(L,L')=a\,\frac{C_\chi(L,L')}{C_{\chi'}(L,L')}e_{\chi'}(L,L').
\]
Given the reference family's two lattice-change laws, this identity transports the target intertwiner with exponent \(-\lfloor N/2\rfloor\) and the source intertwiner with \(\lfloor N'/2\rfloor\): multiply the quotient by the relative determinant factors from each character, and cancel the reference factor using its nonvanishing. Under simultaneous unitary change the determinant-of-basis terms coincide for the two characters and cancel from the quotient. Establishing the reference family's three laws and its comparison at the standard lattices remains open in this note (\nolinkurl{GlobalTheory.Mem.NatEqui.eqP}, \path{IwasawaDecomposition/NormalisedFamilySatisfiesEquivarianceLaws.lean:17--33,42--60,63--117}). The root is therefore \emph{partial}. Its four conclusions enter Z4.2.09 and the class-kernel pullback in Z4.2.10--11. Z4.2.12 inherits them through those two calls.

\emph{Counting the shared input.} Following both type and value edges from \nolinkurl{GlobalTheory.TieBase.base_data} (\path{SatakeParameters/KernelEntryPeriodNeZeroAndLatticeLaws.lean:14--19}) gives 60,733 inclusive core-declaration lines and eight literature-boundary declarations in the project ledger. All 60,733 lines also occur under other expert assignments. Among the 6,438 core theorem declarations, 3,766 have valid ranges. They contribute 51,444 lines with both endpoints included. Omitting the 176 one-line ranges reproduces the independent review's figure of 51,268. Listing the input separately changes the assignment count and leaves the overall union unchanged.

